\documentclass[pdflatex,sn-mathphys-num]{sn-jnl}

\usepackage{amsmath,amssymb,amsfonts}%
\usepackage{amsthm}%
\usepackage{mathtools}%
\usepackage{enumitem}%

\hbadness=10000
\vbadness=10000

%% Theorem environments
\theoremstyle{thmstyleone}%
\newtheorem{theorem}{Theorem}[section]
\newtheorem{proposition}[theorem]{Proposition}
\newtheorem{corollary}[theorem]{Corollary}
\newtheorem{lemma}[theorem]{Lemma}
\newtheorem{conjecture}[theorem]{Conjecture}

\theoremstyle{thmstyletwo}%
\newtheorem{remark}[theorem]{Remark}

\theoremstyle{thmstylethree}%
\newtheorem{definition}[theorem]{Definition}

%% Math operators
\DeclareMathOperator{\tr}{tr}
\DeclareMathOperator{\Hess}{Hess}
\DeclareMathOperator{\Ric}{Ric}
\DeclareMathOperator{\OO}{O}
\DeclareMathOperator{\Sym}{Sym}
\DeclareMathOperator{\Hom}{Hom}

%% Shorthand
\newcommand{\pa}{\partial}
\newcommand{\nn}{\nabla}
\newcommand{\del}{\delta}
\newcommand{\half}{\tfrac{1}{2}}

\numberwithin{equation}{section}

% ================================================================
\begin{document}
% ================================================================

\title[The Einstein Tensor in Seam Geometry]{The Einstein Tensor in Seam Geometry:
  Generator Closure, Emergent Signature, and Null-Convexity}

\author*[1]{\fnm{Lars} \sur{R\"{o}nnb\"{a}ck}}\email{lars@uptochange.com}

\affil*[1]{\orgdiv{Department of Computer and Systems Sciences},
  \orgname{Stockholm University},
  \orgaddress{\city{Stockholm}, \country{Sweden}}}

\abstract{We introduce \emph{seam geometry}, a framework in which a
(pseudo-)Riemannian metric is generated from a single scalar field
through four equivariant tensor operations---conformal scaling,
gradient tensor, Hessian, and Hessian square.  For three maximally
symmetric ansatz classes (conformally flat, spherically symmetric, and
Lorentz-invariant), we compute the Einstein tensor explicitly and show
that it lies in the generator span; in each case the vacuum condition
forces flatness, and metric signature emerges as an output of the
construction.

Identifying the seam scalar with the log-amplitude of a Klein--Gordon
field connects the framework to quantum mechanics: the vacuum condition
becomes the eikonal equation for the phase, and requiring the field to
propagate self-consistently on the metric it generates yields a
nonlinear eigenvalue problem.  A cube-power substitution linearises the
massless case into a Helmholtz equation whose ground-state
positivity constraint uniquely quantises the frequency; the massive
generalisation exhibits a charge-dependent fold producing discrete
weakly and strongly gravitating branches.

Promoting the seam to a complex-valued field with six generators
resolves each structural limitation of the real framework---single
degree of freedom, vanishing vorticity, and vacuum flatness
rigidity---and the Taub plane-symmetric vacuum is realised as an
explicit complex seam metric.  Setting the two real components equal to
the amplitude and phase of a wavefunction defines a self-consistent
system in which a quantum field generates the curved spacetime it
propagates on; the conformal limit recovers the spectrally quantised
system.}

\keywords{seam geometry, Einstein tensor, wavefunction metric,
spectral quantisation, complex seam, vacuum rigidity}

\pacs[MSC Classification]{83C05, 53C50, 53B30, 83C57, 53C21}

\maketitle

% ================================================================
\section{Introduction}
% ================================================================

Suppose one wishes to build a symmetric $(0,2)$-tensor
(a candidate metric) from a scalar field $s$ and a background metric~$h$,
using at most second-order derivatives of $s$ and respecting the isotropy
of~$h$.  On dimensional and symmetry grounds the available building
blocks are:
\begin{itemize}[noitemsep]
\item $h_{ij}$ itself (conformal rescaling),
\item $\pa_i s\,\pa_j s$ (gradient tensor),
\item $\nn^2_{ij} s$ (Hessian),
\item $(\nn^2 s)^2_{ij}:=(\pa_i\pa_k s)\,h^{kl}(\pa_l\pa_j s)$
  (Hessian square).
\end{itemize}
The first three are linear in the derivatives of~$s$ of each order;
the fourth is the simplest \emph{nonlinear} (quadratic-in-Hessian)
building block.  Together they form a four-generator decomposition:
\begin{equation}
  \label{eq:ud}
  g_{ij}
  = \alpha\,h_{ij}
    + \beta\,\pa_i s\,\pa_j s
    + \gamma\,\nn^2_{ij} s
    + \delta\,(\nn^2 s)^2_{ij},
\end{equation}
where $\alpha,\beta,\gamma,\delta$ are smooth functions of $s$ and its
lower-order invariants (e.g.\ $|\nn s|^2_h$, $\Delta_h s$).  We call
this the \emph{four-generator seam decomposition} and refer to the four
terms as the \emph{conformal ($\alpha$),
gradient-tensor ($\beta$), Hessian ($\gamma$), and
Hessian-square ($\delta$) generators}
respectively.  We call~$s$ the \emph{seam}.

\begin{definition}[Four-generator seam decomposition]
\label{def:seam-decomposition}
Let $(M^n,h)$ be a (pseudo-)Riemannian manifold with $n\ge 2$ and
$s\colon M\to\mathbb{R}$ a smooth scalar field.  The
\emph{seam-generated metric} is
\begin{equation}
\label{eq:seam-metric}
g = \alpha\,h + \beta\,ds\otimes ds + \gamma\,\nabla^2 s
    + \delta\,(\nabla^2 s)^2,
\end{equation}
where $\alpha,\beta,\gamma,\delta\colon M\to\mathbb{R}$ are smooth scalar
coefficients and $(\nabla^2 s)^2_{\mu\nu}:=(\nabla_\mu\nabla_\alpha s)\,
h^{\alpha\beta}(\nabla_\beta\nabla_\nu s)$.  By standard invariant theory
for symmetric bilinear forms~\cite{KMS1993}, the five $O(n)$-invariant
scalars of the 2-jet of $s$,
\begin{equation}
  \label{eq:UD-invariants}
  \bigl(s,\; |\nabla s|_h^2,\; \Delta_h s,\; |\nabla^2 s|_h^2,\;
  \langle\nabla s,\nabla^2 s\cdot\nabla s\rangle_h\bigr),
\end{equation}
are available as arguments for $\alpha,\beta,\gamma,\delta$.
\end{definition}

\begin{remark}[Generator collapse in symmetric families]
\label{rem:generator-collapse}
When the Hessian $\nabla^2 s$ is proportional to the background metric
($\nabla^2 s=\lambda\,h$), the Hessian-square generator reduces to
$(\nabla^2 s)^2=\lambda^2 h$, which is absorbed into the conformal
generator~$\alpha$.  This occurs in all three symmetric families studied
in this paper:
\begin{itemize}[noitemsep]
\item \emph{Conformal family} (Result~1): the metric $g=e^{2s}\eta$ uses
  only $\alpha$ (pure conformal rule; $\delta=0$).
\item \emph{$O(n)$-invariant family} (Result~2):
  $\nabla^2 s=\delta_{ij}$ for $s=\tfrac{1}{2}r^2$, so
  $(\nabla^2 s)^2=\delta$.
\item \emph{Lorentz-invariant family} (Result~3):
  $\nabla^2 s=\eta_{\mu\nu}$ for
  $s=\tfrac{1}{2}\eta_{\alpha\beta}x^\alpha x^\beta$, so
  $(\nabla^2 s)^2=\eta$.
\end{itemize}
In these settings the decomposition reduces to three effective generators,
and all computations in Sections~\ref{sec:conformal}--\ref{sec:lorentz}
are carried out in this three-generator form.

Conversely, whenever $\nabla^2 s$ has three or more distinct eigenvalues,
the fourth generator is independent and cannot be absorbed
(Theorem~\ref{thm:closure-fails}).

We note that~\eqref{eq:seam-metric} can also be motivated via the
invariant theory of $O(n)$-representations: the four generators
correspond to the irreducible equivariant maps from $2$-jet data to
symmetric $(0,2)$-tensors (the first three are linear in the Hessian;
the fourth is the simplest quadratic contribution).  This
invariant-theoretic perspective is not needed for any result in this
paper; we mention it only to note that the four-generator span is
natural and not ad hoc.
\end{remark}

\noindent\textbf{Physical motivation.}\enspace
Why study metrics generated entirely by a single scalar?  Three reasons.
First, many physically important spacetimes are conformally flat or have
high symmetry, and in those cases the metric \emph{is} determined by a
scalar to a very good approximation (FLRW, isotropic Schwarzschild
coordinates, cosmological perturbation theory in the Newtonian gauge).
Second, the seam framework provides an ``effective matter'' dictionary
(Section~\ref{sec:discussion}): each generator maps to a specific
stress-energy component (conformal $\to$ isotropic pressure, gradient
$\to$ momentum flux, Hessian $\to$ tidal stress), making the origin of
effective matter transparent.  Third, the self-consistent wavefunction
seam (Section~\ref{subsec:self-consistent}) naturally connects to semiclassical
gravity and Bohmian mechanics, producing spectral quantisation from
geometric self-consistency alone.  The approach is of course
limited---a single scalar cannot encode the full two radiative degrees of
freedom of vacuum GR---but the structural results (closure, rigidity,
quantisation) it delivers are sharp.

\medskip
This paper asks: if the spacetime metric is itself generated by a seam,
what is the Einstein tensor $G_{\mu\nu}$ in terms of these generators?
The individual curvature computations that follow are largely classical
(conformal transformation formulas, warped-product identities, FLRW
dynamics); the contribution here is their systematic organisation within
the seam-generator language and the structural consequences that emerge
from it---closure, rigidity, and the effective-matter dictionary.
We give complete answers for two natural families.

In Section~\ref{sec:conformal} we treat the conformally flat family
$g_{\mu\nu}=e^{2s}\eta_{\mu\nu}$ (pure $\alpha$-rule) and prove that
$G_{\mu\nu}$ is an exact linear combination of all three generators,
with the vacuum condition equivalent to harmonicity of $e^s$.

In Section~\ref{sec:on} we treat the seam decomposition metric generated
from the Euclidean seam $s=\tfrac{1}{2}r^2$, which is the $O(n)$-invariant
family $g_{ij}=A(r)\delta_{ij}+B(r)x_ix_j$.  This family encompasses all
spherically symmetric spatial metrics, including the Schwarzschild and
FLRW 3-slices.  We show that $G_{ij}$ again lies exactly within the span
of the seam generators (now collapsed from four to two by the
degeneracy of the Euclidean Hessian), and that the vacuum condition
forces the metric to be flat within this family---a no-go consequence
of the high symmetry of the scalar seam.

In Section~\ref{sec:lorentz} we treat the Lorentz-invariant seam
$s=\half\eta_{\mu\nu}x^\mu x^\nu$ on Minkowski background.  The resulting
two-function family $g_{\mu\nu}=A(s)\eta_{\mu\nu}+B(s)x_\mu x_\nu$ is
equivalent (inside the light cone) to the general open ($k=-1$) FLRW ansatz.
We compute $G_{\mu\nu}$ explicitly and show that, as in the standard open
($k=-1$) FLRW vacuum setting, vacuum forces flatness within this family.

Section~\ref{sec:signature} discusses the signature-change and a
sphericity consequence of equivariance.

In Section~\ref{sec:wavefunction} we apply the conformal seam framework to
a complex scalar field $\psi=|\psi|\,e^{i\theta}$ by setting
$s=\log|\psi|$.  The seam-geometric vacuum condition acquires a sharp
quantum-mechanical meaning: it is equivalent to the eikonal equation for
the phase when $\psi$ satisfies the Klein--Gordon equation.

In Section~\ref{sec:complex} we promote the seam to a complex-valued
field $\varphi=u+iv$ and show that all three structural
limitations of the real ansatz---vacuum flatness rigidity, a single
degree of freedom, and irrotationality---are resolved by the complex
extension.  Vacuum flatness persists in the conformal-plus-gradient
subfamily, but the full six-generator span admits genuinely curved
vacua (Theorem~\ref{thm:complex-vacuum}).  Setting
$u = \log|\psi|$ and $v = \theta$ connects the complex seam to the
wavefunction framework, defining a self-consistent system in which a
quantum field generates the curved spacetime it propagates on.

\medskip
\noindent\textbf{A no-go / rigidity viewpoint.}
In both symmetry-reduced seam-generated families where a vacuum theorem is
available (Sections~\ref{sec:on} and~\ref{sec:lorentz}), the Einstein equations
do not merely \emph{select} a distinguished metric within the ansatz; they
\emph{collapse} the ansatz to flatness (up to diffeomorphism and scale).
We view this as a no-go statement about these symmetry-reduced ansatz classes:
once the seam rule reduces the geometry to the spherically symmetric spatial
or open-FLRW forms studied here, the remaining vacuum equations leave no room
for nontrivial vacuum curvature within the class.

At the same time, because the metric is generated from a single scalar seam,
the present framework is not expected to encode the full two local radiative
degrees of freedom of vacuum general relativity.  The focus of this paper is
therefore structural: closure, rigidity, signature, and effective-matter
translations within a deliberately restrictive class of scalar-generated
metrics.

\subsection*{Relation to prior work}

The conformally flat Einstein tensor is classical~\cite{Wald1984}.
The $O(n)$-invariant spatial metric and its curvature are also
well-known~\cite{MisnerThorneWheeler1973}.
The contribution of the present paper is the organisation:
in the two explicit seam-generated families computed here, the Einstein
tensor remains in the span of the seam generators that built the metric.
In the $O(n)$-invariant Euclidean seam case, this ``closure'' is symmetry-forced
(every $O(n)$-invariant symmetric 2-tensor is a linear combination of
$\delta_{ij}$ and $x_ix_j$).  In the conformal seam case, it reflects the
rigidity of natural second-order operators on metrics.  Extending this
perspective beyond highly symmetric settings is an open direction.

% ================================================================
\section{Result 1: The Conformal Seam}
\label{sec:conformal}
% ================================================================

Let $\eta_{\mu\nu}$ denote the flat Minkowski (or Euclidean) metric in
$n$ dimensions.  Let $s\in C^\infty$ and set
$g_{\mu\nu}=e^{2s}\eta_{\mu\nu}$.

\begin{proposition}[Ricci tensor, conformal seam]
\label{prop:ricci}
\begin{equation}
  \label{eq:ricci}
  R_{\mu\nu}
  = -(n-2)\bigl[\pa_\mu\pa_\nu s - (\pa_\mu s)(\pa_\nu s)\bigr]
    - \eta_{\mu\nu}\bigl[\Box_\eta s + (n-2)|\nn s|^2_\eta\bigr].
\end{equation}
\end{proposition}

\begin{proof}
Standard conformal transformation formula on a flat background;
see~\cite{Wald1984}, Appendix~D.
\end{proof}

\begin{proposition}[Ricci scalar, conformal seam]
\label{prop:ricciscalar}
\begin{equation}
  \label{eq:ricciscalar}
  R = -e^{-2s}(n-1)\bigl[2\,\Box_\eta s + (n-2)|\nn s|^2_\eta\bigr].
\end{equation}
In $n=4$: $R=-6\,e^{-2s}(\Box_\eta s+|\nn s|^2_\eta)$.
\end{proposition}

\begin{proof}
Contract~\eqref{eq:ricci} with $g^{\mu\nu}=e^{-2s}\eta^{\mu\nu}$ and
collect terms.
\end{proof}

\begin{theorem}[Seam decomposition, conformal seam]
\label{thm:conformal}
In $n=4$,
\begin{equation}
  \label{eq:main}
  \boxed{
  G_{\mu\nu}
  = \underbrace{-2\,\pa_\mu\pa_\nu s}_{\gamma\text{-generator}}
    + \underbrace{2\,(\pa_\mu s)(\pa_\nu s)}_{\beta\text{-generator}}
    + \underbrace{e^{-2s}(2\,\Box_\eta s + |\nn s|^2_\eta)\,g_{\mu\nu}}_{\alpha\text{-generator}}.
  }
\end{equation}
The coefficients $\pm 2=\pm(n-2)|_{n=4}$ are exact.
\end{theorem}

\begin{proof}
$G_{\mu\nu}=R_{\mu\nu}-\half R g_{\mu\nu}$.  Substituting
Propositions~\ref{prop:ricci}--\ref{prop:ricciscalar} with $n=4$:
\begin{align*}
  G_{\mu\nu}
  &= -2\bigl[\pa_\mu\pa_\nu s-(\pa_\mu s)(\pa_\nu s)\bigr]
     - \eta_{\mu\nu}(\Box_\eta s+2|\nn s|^2)
     +3\,\eta_{\mu\nu}(\Box_\eta s+|\nn s|^2).
\end{align*}
Collecting and writing $\eta_{\mu\nu}=e^{-2s}g_{\mu\nu}$
gives~\eqref{eq:main}.
The identification with the three generators of
Definition~\ref{def:seam-decomposition} is immediate:
$\pa_\mu\pa_\nu s$ is the Hessian
generator ($\gamma=1$), $(\pa_\mu s)(\pa_\nu s)$ is the gradient-tensor
generator ($\beta=1$), and $g_{\mu\nu}$ is the conformal generator
($\alpha=1$).
\end{proof}

\begin{remark}[Trace check]
$g^{\mu\nu}G_{\mu\nu}=6e^{-2s}(\Box_\eta s+|\nn s|^2)=-R$.~\checkmark
\end{remark}

\begin{remark}[Dimension dependence]
The coefficients $\pm(n-2)$ vanish at $n=2$, consistent with
$G_{\mu\nu}\equiv 0$ in two dimensions.
\end{remark}

\begin{proposition}[Verification on $H^4$]
\label{prop:h4}
For $s=-\log x_4$ (Poincar\'{e} upper half-space model), $G_{\mu\nu}=3g_{\mu\nu}$,
consistent with $\Ric=-(n-1)g$ at constant curvature $\kappa=-1$.
\end{proposition}

\begin{proof}
$\pa_4 s=-1/x_4$, all other partials zero; $\Box_\eta s=|\nn s|^2=1/x_4^2$.
Direct substitution into~\eqref{eq:main} gives $G_{\mu\nu}=3g_{\mu\nu}$.
\end{proof}

\subsection{Physical interpretation of the three generators}
\label{subsec:physical}

Via $8\pi T_{\mu\nu}=G_{\mu\nu}$, each generator corresponds to a
distinct mode of stress-energy.

\medskip\noindent
\textbf{$\gamma$-generator} $(-2\,\pa_\mu\pa_\nu s)$:~
Anisotropic tidal and viscous stress.  This term is a trace-free symmetric
tensor that distinguishes directions; it vanishes iff $s$ is affine.
It has no counterpart in a perfect fluid.  Geometrically, it measures
how the level sets of $s$ curve in different directions and is nonzero
precisely on the degeneration locus of the seam rule (see Theorem~\ref{thm:signature}).

\medskip\noindent
\textbf{$\beta$-generator} $(+2\,(\pa_\mu s)(\pa_\nu s))$:~
Kinetic momentum flux.  This rank-one tensor has the same structure as
the kinetic stress-energy of a scalar field $\phi$: $T^{(\mathrm{kin})}_{\mu\nu}
=\pa_\mu\phi\,\pa_\nu\phi-\half(\pa_\lambda\phi\,\pa^\lambda\phi)g_{\mu\nu}$
(the trace part is absorbed into the $\alpha$-generator).
For pressureless dust with 4-velocity $u_\mu\propto\pa_\mu s$, the
stress-energy is purely of $\beta$-type.

\medskip\noindent
\textbf{$\alpha$-generator} $(e^{-2s}(2\Box_\eta s+|\nn s|^2)\,g_{\mu\nu})$:~
Isotropic pressure and effective cosmological term.  Stress-energy
proportional to $g_{\mu\nu}$ represents a perfect fluid at rest or a
cosmological constant.  Here the coefficient is dynamical: it depends
on $\Box_\eta s$ and $|\nn s|^2$, so the effective cosmological term is
generated by the seam configuration rather than being fixed input.

\begin{remark}[Information hierarchy]
The three generators carry information at ascending derivative order:
the $\alpha$-generator depends on $s$ and its first derivatives through
$|\nn s|^2$ and $\Box s$; the $\beta$-generator on the gradient $\nn s$;
the $\gamma$-generator on the full Hessian $\pa_\mu\pa_\nu s$.  The richer
the second-order structure of the seam, the more of the stress-energy
is carried by the Hessian generator.
\end{remark}

\begin{corollary}[No light deflection for conformal seams]
\label{cor:no-lensing}
For any conformal seam metric $g_{\mu\nu}=e^{2s}\eta_{\mu\nu}$, the
unparameterised null geodesics of $g$ coincide with those of the flat background
$\eta$.  In particular, in background inertial coordinates the light-ray
trajectories are straight lines (zero deflection angle), although affine
parameterisation and focusing properties may change.
\end{corollary}

\begin{proof}
Conformal transformations preserve the null cone and map null geodesics to null
geodesics as unparameterised curves; see, e.g.,~\cite{Wald1984,Carroll2004}.
\end{proof}

\begin{remark}[Limitation: no gravitational lensing]
\label{rem:no-lensing}
Corollary~\ref{cor:no-lensing} implies that the conformally flat family
cannot describe gravitational lensing, since the standard Schwarzschild
spacetime is not globally conformally flat.  Lensing requires the full
$g_{ij}=A\delta_{ij}+B x_ix_j$ form with a non-trivial gradient-tensor
generator (as can be verified from the isotropic Schwarzschild
decomposition, where the gradient-tensor coefficient vanishes and the
metric lies in the conformal-plus-Hessian span).
\end{remark}

\subsection{Vacuum condition}
\label{subsec:vacuum}

\begin{theorem}[Seam-geometric vacuum condition, conformal case]
\label{thm:vacuum}
$G_{\mu\nu}=0$ if and only if:
\begin{enumerate}[label=(\roman*)]
\item $e^s$ is harmonic in the flat background:~$\Box_\eta e^s=0$;
\item $\pa_\mu\pa_\nu s = (\pa_\mu s)(\pa_\nu s)
  -\half|\nn s|^2\eta_{\mu\nu}$ for all $\mu,\nu$.
\end{enumerate}
\end{theorem}

\begin{proof}
Taking the trace of~\eqref{eq:main}: $G=0$ gives
$\Box_\eta s+|\nn s|^2=0$.  Since
$\Box_\eta(e^s)=e^s(\Box_\eta s+|\nn s|^2)$,
this is equivalent to~(i).  Given~(i), the
$\alpha$-generator evaluates to $-|\nn s|^2\eta_{\mu\nu}$,
and setting the full expression to zero gives~(ii).
\end{proof}

\begin{remark}
By Liouville's theorem, a positive harmonic function on a complete
flat manifold is constant.  Condition~(i) therefore forces $e^s=\mathrm{const}$
(i.e., $g_{\mu\nu}=\mathrm{const}\cdot\eta_{\mu\nu}$) unless the domain
is non-compact or has boundary.
\end{remark}

\begin{corollary}[A harmonic hierarchy in the conformal family]
\label{cor:harmonic-hierarchy}
In the conformal seam family $g_{\mu\nu}=e^{2s}\eta_{\mu\nu}$ in $n=4$:
\begin{enumerate}[label=(\roman*),noitemsep]
\item The scalar curvature vanishes ($R=0$) if and only if $e^s$ is harmonic in
  the flat background: $\Box_\eta(e^s)=0$.
\item The isotropic (conformal) generator in~\eqref{eq:main} vanishes, i.e.
  $(2\Box_\eta s+|\nabla s|_\eta^2)=0$, if and only if $e^{s/2}$ is harmonic:
  $\Box_\eta(e^{s/2})=0$.
\end{enumerate}
\end{corollary}

\begin{proof}
In $n=4$, Proposition~\ref{prop:ricciscalar} gives
$R=-6e^{-2s}(\Box_\eta s+|\nabla s|_\eta^2)$, hence $R=0\iff \Box_\eta s+|\nabla s|_\eta^2=0$.
Since $\Box_\eta(e^{s})=e^{s}(\Box_\eta s+|\nabla s|_\eta^2)$ and $e^{s}>0$,
this is equivalent to $\Box_\eta(e^{s})=0$.

For (ii), expand $\Box_\eta(e^{s/2})=\tfrac12 e^{s/2}(\Box_\eta s+\tfrac12|\nabla s|_\eta^2)$.
Since $e^{s/2}>0$, $\Box_\eta(e^{s/2})=0$ is equivalent to
$\Box_\eta s+\tfrac12|\nabla s|_\eta^2=0$, i.e.
$2\Box_\eta s+|\nabla s|_\eta^2=0$.
\end{proof}

\begin{proposition}[Harmonic ladder exclusivity]
\label{prop:harmonic-ladder}
Let $s\in C^\infty(\mathbb{R}^{1,3})$.  Fix two distinct constants
$k_1\neq k_2$.  If
\begin{equation}
  \label{eq:harmonic-ladder}
  \Box_\eta(e^{k_1 s})=0
  \qquad\text{and}\qquad
  \Box_\eta(e^{k_2 s})=0,
\end{equation}
then $|\nabla s|_\eta^2\equiv 0$ (and hence $\Box_\eta s\equiv 0$).
\end{proposition}

\begin{proof}
For any constant $k$, one has
$\Box_\eta(e^{ks})=k e^{ks}(\Box_\eta s+k|\nabla s|_\eta^2)$.  Since
$e^{ks}\neq 0$, the conditions~\eqref{eq:harmonic-ladder} imply
$\Box_\eta s=-k_1|\nabla s|_\eta^2=-k_2|\nabla s|_\eta^2$.  Subtracting gives
$(k_1-k_2)|\nabla s|_\eta^2=0$, hence $|\nabla s|_\eta^2=0$ and then
$\Box_\eta s=0$.
\end{proof}

\begin{remark}
Proposition~\ref{prop:harmonic-ladder} subsumes the ``dual wave'' observation
of Proposition~\ref{prop:dual-wave} as the special case $(k_1,k_2)=(1,-1)$.
It also shows, for example, that simultaneously requiring
$\Box_\eta(e^{s})=0$ (scalar-flatness in Corollary~\ref{cor:harmonic-hierarchy})
and $\Box_\eta(e^{s/2})=0$ (vanishing isotropic generator) forces a null seam.
\end{remark}

% ================================================================
\section{Result 2: The Full \texorpdfstring{$O(n)$}{O(n)}-Invariant Family}
\label{sec:on}
% ================================================================

We now apply the seam decomposition (with collapsed fourth generator) to the
Euclidean seam $s=\tfrac{1}{2}r^2=\tfrac{1}{2}|x|^2$ on $\mathbb{R}^n$
with flat background $h=\delta$.

\subsection{Metric and generator collapse}

\begin{proposition}[$O(n)$-invariant output from the Euclidean seam]
\label{prop:on}
The five $O(n)$-invariant jet scalars of $s=\half r^2$ are:
$(s,\,|\nn s|^2,\,\Delta s,\,|\nn^2 s|^2,\,\langle\nn s,\nn^2 s\cdot\nn s\rangle)
= (\half r^2,\,r^2,\,n,\,n,\,r^2)$.
Since $\nn^2_{ij}s=\delta_{ij}=h_{ij}$, the Hessian and conformal
generators produce the same tensor; the seam decomposition
reduces to
\begin{equation}
  \label{eq:onmetric}
  g_{ij} = A(r)\,\delta_{ij} + B(r)\,x_ix_j,
\end{equation}
where $A(r):=\alpha(r)+\gamma(r)>0$ and $B(r)=\beta(r)$.  Here $A$ and $B$
may vary with $r$ because $\alpha,\beta,\gamma$ are allowed to depend on the
jet invariants of $s$, and for $s=\half r^2$ these invariants reduce to
functions of $r^2$.
\end{proposition}

In spherical coordinates, writing $P(r):=A(r)+B(r)r^2$,
\begin{equation}
  \label{eq:spherical}
  \mathrm{d}s^2 = P(r)\,\mathrm{d}r^2 + A(r)\,r^2\,\mathrm{d}\Omega^2_{n-1}.
\end{equation}
Equivalently, any $O(n)$-invariant Riemannian metric can be written as
$\mathrm{d}s^2=U(r)\,\mathrm{d}r^2+V(r)\,r^2\,\mathrm{d}\Omega^2_{n-1}$; this
corresponds to~\eqref{eq:spherical} by taking $V=A$ and $U=P$, i.e.
$B=(U-V)/r^2$.  In particular,~\eqref{eq:onmetric} encompasses all rotationally
symmetric metrics in this coordinate chart (sphere, hyperbolic space, and
standard spatial slices such as Schwarzschild and FLRW).

\subsection{Einstein tensor of the \texorpdfstring{$O(n)$}{O(n)}-invariant family}

Introduce the geodesic radial coordinate $\rho$ by
$\mathrm{d}\rho=\sqrt{P(r)}\,\mathrm{d}r$ and the warp factor
\begin{equation}
  \label{eq:warp}
  f(r):=r\,\sqrt{A(r)}.
\end{equation}
Then~\eqref{eq:spherical} takes the warped-product form
$\mathrm{d}s^2=\mathrm{d}\rho^2+f(\rho)^2\,\mathrm{d}\Omega^2_{n-1}$.
The curvature follows by standard warped-product formulae.

\begin{proposition}[Sectional curvatures]
\label{prop:sect}
\begin{align}
  \label{eq:Kra}
  K(r,a) &= -\frac{1}{f}\,\frac{\mathrm{d}^2 f}{\mathrm{d}\rho^2},
  \\[4pt]
  \label{eq:Kab}
  K(a,b) &= \frac{1-\left(\frac{\mathrm{d} f}{\mathrm{d}\rho}\right)^{\!2}}{f^2},
\end{align}
where $a,b$ denote tangential directions.
\end{proposition}

\begin{proof}
Standard warped-product curvature identities.
\end{proof}

\begin{remark}[Flat space check]
If $A\equiv \mathrm{const}$ and $B\equiv 0$ then $P\equiv A$ and $f(\rho)=\sqrt{A}\,\rho$,
so $K(r,a)=0$, $K(a,b)=0$, $G_{ij}=0$.~\checkmark
\end{remark}

\begin{theorem}[Einstein tensor of the $O(n)$-invariant family]
\label{thm:on}
In $n=4$ (three angular directions), the mixed Einstein tensor
components of~\eqref{eq:onmetric} are
\begin{align}
  \label{eq:Grr}
  G^r{}_r &= -3\,\frac{1-\left(\frac{\mathrm{d} f}{\mathrm{d}\rho}\right)^{\!2}}{f^2},
  \\[6pt]
  \label{eq:Gaa}
  G^a{}_a &= 2\,\frac{1}{f}\,\frac{\mathrm{d}^2 f}{\mathrm{d}\rho^2}
  -\frac{1-\left(\frac{\mathrm{d} f}{\mathrm{d}\rho}\right)^{\!2}}{f^2}.
\end{align}
In Cartesian form, $G_{ij}=C_1(r)\,\delta_{ij}+C_2(r)\,x_ix_j$,
where
\begin{equation}
  \label{eq:C1C2}
  C_1 = G^a{}_a\cdot A(r),\qquad
  C_2 = \frac{G^r{}_r\cdot P(r) - G^a{}_a\cdot A(r)}{r^2}.
\end{equation}
\end{theorem}

\begin{proof}
For the warped product $\mathrm{d}\rho^2+f(\rho)^2\,\mathrm{d}\Omega^2_3$,
the Ricci tensor satisfies $R^r{}_r=-3f_{\rho\rho}/f$ and
$R^a{}_a=2(1-f_{\rho}^2)/f^2-f_{\rho\rho}/f$.  Substituting into
$G=R-\half Rg$ yields~\eqref{eq:Grr}--\eqref{eq:Gaa}.
For the Cartesian form: the radial direction $\hat x$ has covariant eigenvalue
$G^r{}_r\cdot P$ and each tangential direction has covariant eigenvalue
$G^a{}_a\cdot A$, yielding~\eqref{eq:C1C2}.
\end{proof}

\begin{remark}[Closure under the seam decomposition]
\label{rem:closure}
$G_{ij}=C_1\delta_{ij}+C_2 x_ix_j$ involves exactly the two tensors
$\delta_{ij}=\nn^2_{ij}s$ (Hessian generator of $s$) and
$x_ix_j=\pa_is\,\pa_js$ (gradient-tensor generator of $s$).
No new tensor structure appears beyond the three seam
generators.  The Einstein tensor of any $O(n)$-invariant seam metric
lies in the span of the seam's own generators.
In this $O(n)$-invariant setting this is also the only possibility by
invariant theory: there are no other independent $O(n)$-invariant
symmetric rank-2 tensors built from $x$.

Equivalently, since for the Euclidean seam $s=\half|x|^2$ one has
$\nabla^2 s=\delta$ and $ds\otimes ds=x\otimes x$, the Cartesian closure can be
read invariantly as
\[
G = C_1(r)\,\nabla^2 s + C_2(r)\,ds\otimes ds.
\]
\end{remark}

\begin{corollary}[Radial scalar curvature]
\label{cor:scalar-radial}
For an $O(n)$-invariant seam metric of the form~\eqref{eq:spherical}, the scalar
curvature $R$ depends only on the radius (equivalently, only on $\rho$ or $r$).
In particular, in the warped-product variables one has
\begin{equation}
  \label{eq:scalar-warp}
  R = -6\,\frac{f_{\rho\rho}}{f} + 6\,\frac{1-f_{\rho}^2}{f^2}
  \qquad (n=4).
\end{equation}
\end{corollary}

\begin{proof}
This follows from the warped-product form $\mathrm{d}\rho^2+f(\rho)^2\mathrm{d}\Omega_3^2$.
\end{proof}

\subsection{Vacuum condition}

\begin{theorem}[Vacuum condition, $O(n)$-invariant family]
\label{thm:vacuumon}
$G_{ij}=0$ if and only if the metric~\eqref{eq:spherical} is flat.  That is,
within this $O(n)$-invariant seam-generated ansatz class, vacuum forces flatness.
\end{theorem}

\begin{proof}
From~\eqref{eq:Grr}, $G^r{}_r=0$ implies $\mathrm{d}f/\mathrm{d}\rho=\pm 1$.
Then~\eqref{eq:Gaa} gives $\mathrm{d}^2 f/\mathrm{d}\rho^2=0$.  Hence
$f(\rho)=c_1\rho+c_0$ with $c_1^2=1$, so after shifting $\rho$ one has
$f(\rho)=\rho$ up to an overall constant scale.  Proposition~\ref{prop:sect} then
gives $K(r,a)=0$ and $K(a,b)=0$, so the metric is flat.
\end{proof}

\begin{corollary}[$\Lambda$-vacuum forces constant curvature in the $O(n)$ family]
\label{cor:lambda-on}
In the setting of Theorem~\ref{thm:on} (so $n=4$), suppose the metric
\eqref{eq:onmetric} satisfies the Einstein-space condition
\begin{equation}
  \label{eq:einstein-space-on}
  G_{ij}=-\Lambda\,g_{ij}
\end{equation}
for a constant $\Lambda\in\mathbb{R}$.  Then the metric has constant sectional
curvature $\kappa=\Lambda/3$.
Equivalently, in warped-product variables it satisfies
\begin{equation}
  \label{eq:f-lambda-ode}
  f_{\rho\rho}=-\frac{\Lambda}{3}f,
  \qquad
  f_{\rho}^2=1-\frac{\Lambda}{3}f^2.
\end{equation}
\end{corollary}

\begin{proof}
Equation~\eqref{eq:einstein-space-on} is equivalent to $G^r{}_r=G^a{}_a=-\Lambda$.
From~\eqref{eq:Grr} we obtain
$1-f_\rho^2=\tfrac{\Lambda}{3}f^2$, which is the second identity in
\eqref{eq:f-lambda-ode}.
Substituting this into~\eqref{eq:Gaa} and setting $G^a{}_a=-\Lambda$ yields
$2f_{\rho\rho}/f-\Lambda/3=-\Lambda$, hence $f_{\rho\rho}=-(\Lambda/3)f$.
Finally, Proposition~\ref{prop:sect} implies
$K(a,b)=(1-f_\rho^2)/f^2=\Lambda/3$ and $K(r,a)=-f_{\rho\rho}/f=\Lambda/3$.
\end{proof}

\begin{remark}[Schwarzschild slice: $n=3$ versus the $n=4$ computation]
Theorem~\ref{thm:on} computes the Einstein tensor of~\eqref{eq:onmetric}
for $n=4$ (three angular directions).  By contrast, a standard spatial
slice of the Schwarzschild \emph{spacetime} is 3-dimensional ($n=3$,
one radial and two angular directions).  One should therefore not insert
3D Schwarzschild data into the $n=4$ formulas~\eqref{eq:Grr}--\eqref{eq:Gaa}.

For the Schwarzschild spatial 3-metric written in areal radius,
\[
\mathrm{d}s^2 = \frac{r}{r-2M}\,\mathrm{d}r^2 + r^2\,\mathrm{d}\Omega_2^2,
\]
one has $A(r)=1$, $P(r)=\frac{r}{r-2M}$ and $B(r)=\frac{P-A}{r^2}=\frac{2M}{r^2(r-2M)}$.
The $n=3$ warped-product formulae then give
\[
G^r{}_r = -\frac{B(r)}{A(r)\,P(r)}=-\frac{2M}{r^3},
\qquad
G^a{}_a = -\frac{P'(r)}{2r\,P(r)^2}=\frac{M}{r^3}.
\]
Thus the Schwarzschild spatial 3-geometry is not vacuum in the intrinsic
3D sense ($G^{(3)}_{ij}\neq 0$), even though it arises as Cauchy data for a
vacuum 4D spacetime.  In fact this slice is time-symmetric ($K_{ij}=0$), so
the vacuum Hamiltonian constraint yields $R^{(3)}=0$
(see, e.g.,~\cite{Wald1984}, Chapter~10); consequently
$G^{(3)}_{ij}=R^{(3)}_{ij}$ is nonzero because the 3-geometry is scalar-flat
but not Ricci-flat.  Theorem~\ref{thm:vacuumon} is therefore a statement about
vacuum \emph{Riemannian} geometry within the seam family: in the $O(n)$-invariant
seam-generated class, the vacuum condition forces the metric to be flat.
\end{remark}

% ================================================================
\section{Result 3: The Lorentz-Invariant Family}
\label{sec:lorentz}
% ================================================================

We now switch from the Euclidean seam $s=\half r^2$ to the Lorentz-invariant
interval seam on Minkowski background.

\subsection{Metric ansatz and FLRW normal form}

Let $\eta_{\mu\nu}=\mathrm{diag}(-1,1,1,1)$ and define
\begin{equation}
  \label{eq:lorentz-seam-main}
  s=\half\,\eta_{\mu\nu}x^\mu x^\nu=\half(-t^2+r^2).
\end{equation}
Then $\pa_\mu\pa_\nu s=\eta_{\mu\nu}$, so the conformal and Hessian generators
collapse.  Any $SO(1,3)$-invariant seam output therefore has the two-function
form
\begin{equation}
  \label{eq:lorentz-inv-metric-main}
  g_{\mu\nu}=A(s)\,\eta_{\mu\nu}+B(s)\,x_\mu x_\nu,
\end{equation}
where $x_\mu:=\eta_{\mu\nu}x^\nu$.

Inside the future light cone $t^2>r^2$, introduce hyperbolic coordinates
$t=\tau\cosh\chi$, $r=\tau\sinh\chi$ (so $\tau>0$ and $s=-\half\tau^2$).  Then
$\eta=-\mathrm{d}\tau^2+\tau^2\mathrm{d}H_3^2$ and
$x_\mu\mathrm{d}x^\mu=-\tau\,\mathrm{d}\tau$, hence
\begin{equation}
  \label{eq:open-flrw-form-main}
  \mathrm{d}s^2
  =-N(\tau)^2\,\mathrm{d}\tau^2+a(\tau)^2\,\mathrm{d}H_3^2,
\end{equation}
with
\begin{equation}
  \label{eq:A-B-to-a-N-main}
  a(\tau)^2=A(\tau)\,\tau^2,\qquad
  N(\tau)^2=A(\tau)-B(\tau)\,\tau^2.
\end{equation}
Thus the Lorentz-invariant seam parameterises the general open ($k=-1$)
FLRW ansatz class.

\subsection{Einstein tensor and vacuum}

\begin{theorem}[Einstein tensor of the Lorentz-invariant family]
\label{thm:lorentz}
In $n=4$, for the metric~\eqref{eq:open-flrw-form-main} with $k=-1$, the mixed
Einstein tensor components are
\begin{align}
  \label{eq:lorentz-G00}
  G^{\tau}{}_{\tau}
  &= -3\left[\left(\frac{1}{N}\frac{a'}{a}\right)^{\!2}-\frac{1}{a^2}\right],
  \\
  \label{eq:lorentz-Gaa}
  G^{a}{}_{a}
  &= -\left[\frac{2}{N}\frac{\mathrm{d}}{\mathrm{d}\tau}\left(\frac{1}{N}\frac{a'}{a}\right)
    +3\left(\frac{1}{N}\frac{a'}{a}\right)^{\!2}-\frac{1}{a^2}\right],
\end{align}
where $a$ is any tangential (hyperbolic) direction and $'$ denotes
$\mathrm{d}/\mathrm{d}\tau$.

Moreover, $G_{\mu\nu}$ is $SO(1,3)$-invariant and therefore has Cartesian form
$G_{\mu\nu}=C_1(s)\,\eta_{\mu\nu}+C_2(s)\,x_\mu x_\nu$ for scalar functions
$C_1,C_2$.

In terms of the FLRW mixed components, these coefficients can be recovered as
\begin{equation}
  \label{eq:lorentz-C1C2}
  C_1 = A\,G^{a}{}_{a},
  \qquad
  C_2 = \frac{-N^2\,G^{\tau}{}_{\tau}+A\,G^{a}{}_{a}}{\tau^2}
  = \frac{-N^2\,G^{\tau}{}_{\tau}+A\,G^{a}{}_{a}}{-2s},
\end{equation}
where $a$ denotes any tangential hyperbolic direction.
\end{theorem}

\begin{proof}
The metric~\eqref{eq:open-flrw-form-main} is a standard FLRW form with spatial
curvature $k=-1$.  Equations~\eqref{eq:lorentz-G00}--\eqref{eq:lorentz-Gaa} are
the classical Einstein tensor identities for FLRW in lapse form; see, e.g.,~\cite{Weinberg1972,Carroll2004}.
The final statement follows because the only $SO(1,3)$-invariant covariant
rank-2 tensors built from $x$ and $\eta$ are $\eta_{\mu\nu}$ and $x_\mu x_\nu$.
\end{proof}

\begin{corollary}[Vacuum structure of the Lorentz-invariant family]
\label{cor:lorentz-vacuum}
If $G_{\mu\nu}=0$ for a Lorentz-invariant seam metric
\eqref{eq:lorentz-inv-metric-main} in a region inside the light cone, then the
metric is flat (isometric to a constant multiple of Minkowski space).  In other
words, within this Lorentz-invariant seam-generated ansatz class, vacuum forces
flatness.
\end{corollary}

\begin{proof}
In cosmic time gauge $N\equiv 1$, the vacuum equation $G^{\tau}{}_{\tau}=0$ gives
$(\dot a/a)^2=1/a^2$, i.e., $\dot a=\pm 1$.  Hence $a(t)=t+c$ (up to time
reversal), which is the \emph{Milne universe}: an open ($k=-1$) FLRW chart on
flat Minkowski spacetime (the interior of a light cone).
Allowing an arbitrary lapse corresponds only to reparameterisation of time.
Finally, an overall constant scale factor multiplies $\eta$ and preserves
flatness.
\end{proof}

\begin{corollary}[$\Lambda$-vacuum in the Lorentz-invariant family]
\label{cor:lorentz-lambda}
In the Lorentz-invariant seam family, the Einstein-space condition
\begin{equation}
  \label{eq:lorentz-einstein-space}
  G_{\mu\nu}=-\Lambda\,g_{\mu\nu}
\end{equation}
is equivalent (inside the light cone) to the standard open ($k=-1$) FLRW
$\Lambda$-vacuum equations for the scale factor $a$ in
\eqref{eq:open-flrw-form-main}.  In particular, in cosmic time gauge $N\equiv 1$
these reduce to
\begin{equation}
  \label{eq:lorentz-lambda-friedmann}
  \left(\frac{\dot a}{a}\right)^{\!2}=\frac{1}{a^2}+\frac{\Lambda}{3},
  \qquad
  \frac{\ddot a}{a}=\frac{\Lambda}{3}.
\end{equation}
For $\Lambda=0$ this recovers the Milne solution of
Corollary~\ref{cor:lorentz-vacuum}; for $\Lambda\neq 0$ the solutions are the
usual constant-curvature ($\Ric=\Lambda g$) FLRW spacetimes in open slicing.
\end{corollary}

\begin{proof}
Equation~\eqref{eq:lorentz-einstein-space} is equivalent to the two conditions
$G^{\tau}{}_{\tau}=G^{a}{}_{a}=-\Lambda$.  Using
\eqref{eq:lorentz-G00}--\eqref{eq:lorentz-Gaa} with $k=-1$ yields the standard
Friedmann--acceleration system; in cosmic time gauge $N\equiv 1$ this is exactly
\eqref{eq:lorentz-lambda-friedmann}.
\end{proof}

\begin{theorem}[Signature and degeneration (Lorentz-invariant family)]
\label{thm:lorentz-signature}
Inside the light cone ($s<0$), the Lorentz-invariant metric
\eqref{eq:lorentz-inv-metric-main} has the FLRW form~\eqref{eq:open-flrw-form-main}
with $a(\tau)^2=A(\tau)\tau^2$ and $N(\tau)^2=A(\tau)-B(\tau)\tau^2$.
Assuming $A>0$ (so the spatial part is positive-definite), the metric is:
\begin{enumerate}[label=(\roman*)]
\item Lorentzian iff $N(\tau)^2>0$, equivalently $A(s)+2s\,B(s)>0$;
\item degenerate iff $N(\tau)^2=0$, equivalently $A(s)+2s\,B(s)=0$;
\item Riemannian iff $N(\tau)^2<0$.
\end{enumerate}
The degeneration locus $\{N=0\}$ is a level set of $s$, hence a two-sheeted
hyperboloid $t^2-r^2=\tau_0^2$ (or, equivalently, $s=s_0=-\half\tau_0^2$).
\end{theorem}

\begin{proof}
In~\eqref{eq:open-flrw-form-main} one has $g_{\tau\tau}=-N(\tau)^2$ and spatial
metric $a(\tau)^2\,\mathrm{d}H_3^2$.  With $A>0$, the spatial part is positive.
The sign of $N^2$ therefore determines whether the $\tau$-direction is timelike
(Lorentzian), null (degenerate), or spacelike (Riemannian).  Since
$s=-\half\tau^2$, the condition $N(\tau)^2=0$ is equivalent to
$A(s)+2sB(s)=0$, and the level set $\{\tau=\tau_0\}$ is the hyperboloid
$t^2-r^2=\tau_0^2$.
\end{proof}

\begin{remark}[Physical meaning of the Lorentz seam: Weyl's Postulate]
\label{rem:weyl-postulate}
Inside the future light cone, $s=-\half\tau^2$ where
$\tau=\sqrt{t^2-r^2}$ is the proper time along radial geodesics from the
origin.  The gradient $\nabla_\mu s=x_\mu$ is timelike, and the normalised
4-velocity $U_\mu=x_\mu/\tau$ is exactly the comoving Hubble-flow velocity of
the open FLRW model.  By Corollary~\ref{cor:irrotational}, this flow is
irrotational and hypersurface-orthogonal, with the constant-$\tau$ hyperboloids
as spatial slices.  The Lorentz-invariant seam therefore realises
\emph{Weyl's Postulate}---the assumption that the cosmic substratum consists of
a congruence of non-intersecting, irrotational geodesic observers diverging
from a single past event---as a consequence of $SO(1,3)$-equivariance of the
scalar seam, rather than as an independent axiom.
\end{remark}

% ================================================================
\section{Signature as Output and a Sphericity Corollary}
\label{sec:signature}
% ================================================================

\subsection{Signature change at the degeneration sphere}

\begin{theorem}[Signature change]
\label{thm:signature}
The metric~\eqref{eq:onmetric} is:
\begin{enumerate}[label=(\roman*)]
\item positive-definite (Riemannian) at $x$ iff
  $A(r)>0$ and $P(r)=A(r)+B(r)r^2>0$;
\item degenerate at radius $r_0$ iff $A(r_0)=0$ or $P(r_0)=0$; in particular,
  if $A(r)>0$ for all $r$ then the degeneration locus is the sphere $|x|=r_0$
  where $P(r_0)=0$;
\item Lorentzian (signature $(1,n-1)$) at $x$ iff $P(r)<0$.
\end{enumerate}
Metric signature is an output of $(A,B)$, not an independent postulate.
\end{theorem}

\begin{proof}
The eigenvalues of $g_{ij}=A\delta_{ij}+Bx_ix_j$ are
$A$ (on $x^\perp$, multiplicity $n-1$) and $P=A+Br^2$ (on $\hat x$).
The degeneracy locus $\{P=0\}$ is a sphere by the $O(n)$-invariance of $P(r)$.
\end{proof}

\subsection{Sphericity from equivariance}

\begin{corollary}[Sphericity via seam equivariance]
\label{cor:birkhoff}
For any metric generated from the Euclidean seam via the
seam decomposition, every degeneration locus is a sphere.
In particular, every horizon and every signature-change surface
in this family is spherical.
\end{corollary}

\begin{proof}
$P(r)$ depends only on $r=|x|$ by $O(n)$-invariance of the seam
rule.  Any locus $\{P=0\}$ is therefore $O(n)$-invariant, hence a sphere.
\end{proof}

\begin{remark}[Relation to the classical Birkhoff theorem]
The classical Birkhoff theorem states that every spherically symmetric
vacuum solution of the Einstein equations is isometric to a portion of
the Schwarzschild solution~\cite{Birkhoff1923}.  Corollary~\ref{cor:birkhoff}
establishes only the \emph{sphericity} of the relevant loci by $O(n)$-equivariance
of the seam rule, without invoking the Einstein equations.  It does not assert
the dynamical/staticity content of the classical theorem, nor the Schwarzschild
uniqueness conclusion.  Rather, it isolates the prior structural statement that
any degeneration locus in this $O(n)$-generated family must be spherical.
\end{remark}

% ================================================================
\section{Discussion}
\label{sec:discussion}
% ================================================================

\subsection{Generator closure under \texorpdfstring{$G_{\mu\nu}$}{Gmunu}}

The seam decomposition (Definition~\ref{def:seam-decomposition}) employs
four generators.  In each of Results~1--3, the Hessian-square generator
collapses into the span of the other three
(Remark~\ref{rem:generator-collapse}), and the Einstein tensor of the
resulting metric is again expressible in seam variables.
In the $O(n)$-invariant Euclidean seam case this collapse is
symmetry-forced; in the conformal seam case, only the $\alpha$-generator
contributes, so the fourth generator is absent entirely.

For a general seam, however, the Hessian-square generator is
independent once $\nn^2 s$ has three or more distinct eigenvalues.
Below we show that with all four generators present, first-order
closure holds, while any attempt to close with only three generators
fails in general.

\subsubsection*{Perturbative analysis: the Hessian-squared obstruction}

To test closure beyond symmetry we compute the linearised Einstein tensor
of a general seam perturbation on flat background.  Write
$g_{\mu\nu}(\varepsilon)=\eta_{\mu\nu}+\varepsilon h_{\mu\nu}$ with
$h_{\mu\nu}=\alpha_1\,\eta_{\mu\nu}+\beta_1\,\pa_\mu s\,\pa_\nu s
+\gamma_1\,\pa_\mu\pa_\nu s$
(constants $\alpha_1,\beta_1,\gamma_1$).  Each generator contributes
independently at first order.

\begin{lemma}[Linearised $G$ for the gradient-tensor generator]
\label{lem:lin-beta}
On flat $(\mathbb{R}^n,\eta)$, the linearised Ricci tensor of the
perturbation $h_{\mu\nu}=\pa_\mu s\,\pa_\nu s$ is
\begin{equation}\label{eq:lin-ricci-beta}
  R^{(1)}_{\mu\nu}
  = (\Box_\eta s)\,\pa_\mu\pa_\nu s
    - (\nn^2_\eta s)^2_{\mu\nu},
\end{equation}
where $(\nn^2_\eta s)^2_{\mu\nu}
:=(\pa_\mu\pa_\alpha s)\,\eta^{\alpha\beta}(\pa_\beta\pa_\nu s)$
is the Hessian square---the fourth seam generator
(Definition~\ref{def:seam-decomposition}).
The linearised Einstein tensor is
\begin{equation}\label{eq:lin-einstein-beta}
  \boxed{
  G^{(1)}_{\mu\nu}
  = (\Box_\eta s)\,\pa_\mu\pa_\nu s
    - (\nn^2_\eta s)^2_{\mu\nu}
    - \tfrac{1}{2}\bigl[(\Box_\eta s)^2 - |\nn^2_\eta s|^2\bigr]\,
      \eta_{\mu\nu}.
  }
\end{equation}
\end{lemma}

\begin{proof}
The standard linearised Ricci formula on flat background reads
$2R^{(1)}_{\mu\nu}
=\pa_\alpha\pa_\mu h^{\alpha}{}_\nu
+\pa_\alpha\pa_\nu h^{\alpha}{}_\mu
-\Box_\eta h_{\mu\nu}
-\pa_\mu\pa_\nu h$,
where indices are raised and $h:=\eta^{\mu\nu}h_{\mu\nu}$.
For $h_{\mu\nu}=\pa_\mu s\,\pa_\nu s$ one computes
$\pa_\alpha h^{\alpha}{}_\nu
=(\Box_\eta s)\,\pa_\nu s+(\nn^2_\eta s\cdot\nn_\eta s)_\nu$.
Taking one more derivative:
\begin{align*}
  \pa_\mu\pa_\alpha h^{\alpha}{}_\nu
  &= (\pa_\mu\Box_\eta s)\,\pa_\nu s
     + (\Box_\eta s)\,\pa_\mu\pa_\nu s
     + (\nn^2_\eta s)^2_{\mu\nu}
     + (\pa^\alpha s)\,\pa_\mu\pa_\alpha\pa_\nu s.
\end{align*}
Likewise,
$\Box_\eta h_{\mu\nu}
= (\pa_\mu\Box_\eta s)\,\pa_\nu s
  + \pa_\mu s\,(\pa_\nu\Box_\eta s)
  + 2(\nn^2_\eta s)^2_{\mu\nu}$
and
$\pa_\mu\pa_\nu h
= \pa_\mu\pa_\nu|\nn_\eta s|^2
= 2(\nn^2_\eta s)^2_{\mu\nu}
  + 2(\pa^\alpha s)\,\pa_\mu\pa_\nu\pa_\alpha s$.
All third-derivative and $\pa\Box s$ terms cancel, leaving
$2R^{(1)}_{\mu\nu}
= 2(\Box_\eta s)\,\pa_\mu\pa_\nu s - 2(\nn^2_\eta s)^2_{\mu\nu}$.
Tracing gives $R^{(1)}=(\Box_\eta s)^2-|\nn^2_\eta s|^2$, and
$G^{(1)}=R^{(1)}-\half R^{(1)}\eta$ yields~\eqref{eq:lin-einstein-beta}.
\end{proof}

\begin{lemma}[Linearised $G$ for the conformal and Hessian generators]
\label{lem:lin-alpha-gamma}
On flat background:
\begin{enumerate}[label=(\roman*)]
\item For $h_{\mu\nu}=a(s)\,\eta_{\mu\nu}$ (conformal generator with
  $a$ depending on $s$),
\begin{equation}\label{eq:lin-conformal}
  G^{(1)}_{\mu\nu}
  = \tfrac{n{-}2}{2}\bigl[a''|\nn_\eta s|^2+a'\Box_\eta s\bigr]\eta_{\mu\nu}
    + \tfrac{2{-}n}{2}a''\,\pa_\mu s\,\pa_\nu s
    + \tfrac{2{-}n}{2}a'\,\pa_\mu\pa_\nu s.
\end{equation}
Only the three original generators appear.
\item For $h_{\mu\nu}=\gamma_1\,\pa_\mu\pa_\nu s$ (constant $\gamma_1$),
  $G^{(1)}_{\mu\nu}\equiv 0$.  (This perturbation is a pure flat-background
  diffeomorphism gauge mode: $\pa_\mu\pa_\nu s=\half\mathcal{L}_{\nn s}\eta_{\mu\nu}$.)
\end{enumerate}
\end{lemma}

\begin{proof}
Part~(i) follows by direct substitution into the linearised Ricci formula;
the chain-rule derivatives of $a(s)$ introduce only $\pa_\mu s$ and
$\pa_\mu\pa_\nu s$ factors.
Part~(ii): on flat background, partial derivatives commute, so
$\pa_\alpha\pa_\mu h^{\alpha}{}_\nu
=\gamma_1\,\pa_\mu\pa_\nu(\Box_\eta s)
=\gamma_1\,\Box_\eta(\pa_\mu\pa_\nu s)
=\Box_\eta h_{\mu\nu}$, and all terms in $R^{(1)}_{\mu\nu}$ cancel.
\end{proof}

\begin{theorem}[Necessity of the fourth generator]
\label{thm:closure-fails}
Let $n\ge 3$, $\eta$ the flat metric on $\mathbb{R}^n$, and
$s\in C^\infty(\mathbb{R}^n)$.  Consider the seam-generated metric
$g_{\mu\nu}(\varepsilon)=\eta_{\mu\nu}+\varepsilon\,\pa_\mu s\,\pa_\nu s$
for small $\varepsilon>0$.  If
$\nn^2_\eta s$ has at least three distinct eigenvalues at a point $x_0$
where $\nn_\eta s(x_0)=0$, then the first-order Einstein tensor
$G^{(1)}_{\mu\nu}$ does not lie in
$\mathrm{span}\{\eta_{\mu\nu},\,\pa_\mu s\,\pa_\nu s,\,\pa_\mu\pa_\nu s\}$
near $x_0$; the Hessian-square generator is
essential.
\end{theorem}

\begin{proof}
At any point where $\nn s=0$, the gradient-tensor generator vanishes, so the
three-generator span reduces to
$\mathrm{span}\{\eta_{\mu\nu},\,\pa_\mu\pa_\nu s\}$.
By~\eqref{eq:lin-einstein-beta}, membership requires
\[
(\nn^2_\eta s)^2_{\mu\nu}
= c_1\,\eta_{\mu\nu}+c_2\,\pa_\mu\pa_\nu s
\]
for some scalars $c_1,c_2$.  In a basis that diagonalises $\nn^2_\eta s$
with eigenvalues $\lambda_1,\ldots,\lambda_n$, this reads
$\lambda_i^2=c_1+c_2\lambda_i$ for every $i$.  This is a single quadratic
in $\lambda_i$, hence admits at most two solutions.  Three or more distinct
eigenvalues overdetermine $(c_1,c_2)$.

Concretely, take $s=x_1^2+2x_2^2$ on $\mathbb{R}^3$ with $x_0=0$.
Then $\nn^2 s=\mathrm{diag}(2,4,0)$,
$(\nn^2 s)^2=\mathrm{diag}(4,16,0)$,
$\Box s=6$, $|\nn^2 s|^2=20$, and~\eqref{eq:lin-einstein-beta} gives
$G^{(1)}|_0=\mathrm{diag}(0,0,-8)$.
Requiring $\mathrm{diag}(0,0,-8)=c_1\,\mathrm{diag}(1,1,1)+c_2\,\mathrm{diag}(2,4,0)$
gives $c_1+2c_2=0$, $c_1+4c_2=0$, $c_1=-8$:
the first two equations force $c_2=0$ and $c_1=0$, contradicting $c_1=-8$.
\end{proof}

\begin{remark}[Why three generators suffice in Results~1--3]
\label{rem:why-closure-held}
In each of the three computed families, a structural feature renders
the fourth generator redundant:
\begin{enumerate}[label=(\roman*)]
\item \emph{Conformal family} (Result~1):
  $g=e^{2s}\eta$ is a pure $\alpha$-rule.
  By Lemma~\ref{lem:lin-alpha-gamma}(i), only three-generator terms appear
  (the $\beta$-coefficient is zero, so the Hessian-squared obstruction in
  Lemma~\ref{lem:lin-beta} is absent).
\item \emph{$O(n)$-invariant family} (Result~2):
  $\nn^2_\eta s=\delta_{ij}$ for $s=\half r^2$,
  so $(\nn^2 s)^2=\delta$, which is already in the conformal generator span.
\item \emph{Lorentz-invariant family} (Result~3):
  $\nn^2_\eta s=\eta_{\mu\nu}$ for $s=\half\eta_{\alpha\beta}x^\alpha x^\beta$,
  so $(\nn^2 s)^2=\eta$, again collapsing into the span.
\end{enumerate}
In cases~(ii)--(iii) the Hessian has a single eigenvalue (it is
proportional to the background), so $(\nn^2 s)^2$ collapses into
the three-generator span and the fourth generator is automatically
absorbed.
\end{remark}

\begin{theorem}[First-order closure of the seam decomposition]
\label{thm:extended-closure}
At first order in $\varepsilon$, the Einstein tensor of any seam perturbation
$g(\varepsilon)=\eta+\varepsilon(\alpha_1\,\eta+\beta_1\,ds\otimes ds
+\gamma_1\,\nn^2 s)$ (with constant coefficients) lies in the span of
the four seam generators:
\begin{equation}\label{eq:extended-closure}
  G^{(1)}_{\mu\nu}
  \in\mathrm{span}\bigl\{
    \eta_{\mu\nu},\;\;
    \pa_\mu s\,\pa_\nu s,\;\;
    \pa_\mu\pa_\nu s,\;\;
    (\nn^2_\eta s)^2_{\mu\nu}
  \bigr\},
\end{equation}
with scalar coefficients that are jet invariants of $s$.  The fourth
generator $(\nn^2_\eta s)^2$ arises solely from the gradient-tensor
generator (coefficient $\beta_1$) and enters with coefficient $-\beta_1$.
For $\beta_1=0$ or whenever $\nn^2_\eta s\propto\eta$,
the span reduces to the first three generators.
\end{theorem}

\begin{proof}
Combine Lemmas~\ref{lem:lin-beta} and~\ref{lem:lin-alpha-gamma}.
The conformal generator contributes only three-generator
terms~\eqref{eq:lin-conformal}.  The Hessian generator contributes
nothing (gauge mode).  The gradient-tensor generator contributes
$\beta_1$ times~\eqref{eq:lin-einstein-beta}, which introduces
$(\nn^2 s)^2$ with coefficient $-\beta_1$.  The fourth generator is
absent precisely when $\beta_1=0$ or when $\nn^2 s=\lambda\,\eta$ for some
scalar $\lambda$ (in which case
$(\nn^2 s)^2=\lambda^2\eta\in\mathrm{span}\{\eta\}$).
\end{proof}

\noindent\textbf{The closure picture.}\enspace
Theorem~\ref{thm:closure-fails} shows that dropping the fourth
generator breaks closure once the Hessian of the seam has three or more
distinct eigenvalues.
Theorem~\ref{thm:extended-closure} confirms that retaining all four
generators of the seam decomposition
(Definition~\ref{def:seam-decomposition}) restores first-order closure.
The fourth generator $(\nn^2 s)^2$ is the simplest quadratic
$O(n)$-equivariant built from the Hessian; a full classification of the
nonlinear equivariant ring---including the mixed equivariant
$ds\odot(\nn^2 s\cdot ds)$---and the question of whether closure at
the nonlinear level holds within such an extended ring, remain open.

\begin{conjecture}[Nonlinear closure for the conformal--gradient sub-family]
\label{conj:nonlinear-closure}
Let $g=\alpha(s)\,\eta+\beta(s)\,ds\otimes ds$ be a seam metric with
$\gamma=\delta=0$ on flat background $(\mathbb{R}^n,\eta)$.  Then the
full (non-perturbative) Einstein tensor $G_{\mu\nu}(g)$ lies in the
four-generator span
\[
  G_{\mu\nu}
  \in\mathrm{span}\bigl\{
    \eta_{\mu\nu},\;
    \pa_\mu s\,\pa_\nu s,\;
    \pa_\mu\pa_\nu s,\;
    (\nn^2_\eta s)^2_{\mu\nu}
  \bigr\}
\]
with coefficients that are smooth functions of the $2$-jet invariants
of~$s$.
\end{conjecture}

\noindent
\emph{Evidence.}  Theorem~\ref{thm:extended-closure} establishes the
conjecture at first order.  Results~1--3 verify it nonlinearly in three
symmetric sub-families (where the fourth generator collapses).  The
general nonlinear computation for the conformal--gradient family
$g=\alpha\,\eta+\beta\,ds\otimes ds$ involves Christoffel symbols that
are quadratic in the $2$-jet of~$s$; a direct (if lengthy) calculation
should settle the conjecture in this case.

\begin{proposition}[Symmetric closure as an invariant-span statement]
\label{prop:symmetric-closure}
Let $g$ be an $O(n)$-invariant metric on $\mathbb{R}^n$ generated from the
Euclidean seam, so that $g$ has the form~\eqref{eq:onmetric} with radial
coefficients.  Then any $O(n)$-invariant symmetric rank-2 tensor field (in
particular $G(g)$) is necessarily of the form
\[
T_{ij}=C_1(r)\,\delta_{ij}+C_2(r)\,x_ix_j.
\]
\end{proposition}

\noindent\emph{Interpretation:} in this symmetric sector, the Einstein operator acts
as an endomorphism of the two-generator $O(n)$-invariant module spanned by
$\delta$ and $x\otimes x$.

\begin{proof}
At each point with $r\neq 0$, the stabiliser of $x$ is $O(n-1)$, and the space
of $O(n-1)$-invariant symmetric bilinear forms on $\mathbb{R}^n$ is two-dimensional:
one component on $x^\perp$ (proportional to $\delta$) and one on $\mathrm{span}\{x\}$
(proportional to $x\otimes x$).  The coefficients must be radial by $O(n)$-invariance.
\end{proof}

\begin{remark}[Boundary of the symmetric closure]
The span in Proposition~\ref{prop:symmetric-closure} is a symmetry statement: it
captures exactly the tensor structures allowed by $O(n)$-invariance.  For a
generic anisotropic perturbation, additional tensor structures appear.  For
example, the weakly anisotropic conformal metric
\[
g_{ij}=\bigl(1+\varepsilon x_1^2\bigr)\,\delta_{ij}
\]
is not rotationally invariant, and neither is its Einstein tensor.  In
particular, it cannot in general be written as
$C_1(r)\,\delta_{ij}+C_2(r)\,x_ix_j$ with radial coefficients.
\end{remark}

\noindent Results~1--3 provide evidence in three maximally symmetric settings
that the seam-generator viewpoint organises both the Einstein tensor and
the vacuum conditions within scalar-generated ansatz classes.

\subsection{Matter content from generators}

Theorems~\ref{thm:conformal} and~\ref{thm:on} together say that
in seam language, all matter content is one of three things:
anisotropic tidal stress (Hessian generator), kinetic scalar-field
momentum flux (gradient-tensor generator), or isotropic pressure
(conformal generator).  A theory of matter written in seam variables
has a natural three-way decomposition of $T_{\mu\nu}$ built in from
the start.

\begin{remark}[Effective seam matter]
One may regard a seam-generated metric $g=g(s)$ as defining an \emph{effective}
stress--energy tensor by
\begin{equation}
  \label{eq:Teff}
  8\pi\,T^{\mathrm{eff}}_{\mu\nu} := G_{\mu\nu}(g(s)).
\end{equation}
This viewpoint turns seam rules into a constrained ansatz class for
stress--energy models (``seam fluids'') whose algebraic structure is locked to
the jet data of $s$.  It is an effective description: it does not, by itself,
specify a fundamental matter Lagrangian, nor does it imply that any generator
subterm is separately conserved.

Nevertheless, $T^{\mathrm{eff}}$ is \emph{automatically} covariantly conserved
with respect to $g$:
\begin{equation}
  \label{eq:Teff-conserved}
  \nabla^{\mu}T^{\mathrm{eff}}_{\mu\nu}\equiv 0,
\end{equation}
by the contracted Bianchi identity $\nabla^{\mu}G_{\mu\nu}\equiv 0$.
This should be read as a differential identity constraining the seam-generator
coefficients once $G_{\mu\nu}$ is expressed in generator form; it does not
promote any individual generator contribution to a separately conserved fluid.
\end{remark}

\begin{remark}[Imperfect-fluid form (algebraic decomposition)]
Whenever $T^{\mathrm{eff}}_{\mu\nu}$ is evaluated in a region where a preferred
unit timelike vector field $u$ is specified (for example, if $\nabla s$ is
timelike and one sets $u_\mu\propto\pa_\mu s$), one may decompose
$T^{\mathrm{eff}}$ algebraically into energy density, isotropic pressure, and
anisotropic stress in the standard way:
$T^{\mathrm{eff}}_{\mu\nu}=\rho\,u_\mu u_\nu + p\,h_{\mu\nu}+\pi_{\mu\nu}$ with
$h_{\mu\nu}=g_{\mu\nu}+u_\mu u_\nu$ and $\pi$ trace-free and spatial.
This decomposition is kinematical; in general it does not split
$T^{\mathrm{eff}}$ into separately conserved components.
\end{remark}

Energy conditions admit a particularly clean reformulation in seam variables
for the conformal family; see Section~\ref{sec:energy}.

\begin{corollary}[No anisotropic stress without the Hessian generator]
\label{cor:no-anisotropy-no-hess}
Let $g=g(s)$ be a seam-generated metric and define
$T^{\mathrm{eff}}_{\mu\nu}:=G_{\mu\nu}(g)/(8\pi)$.  In any region where
$\nabla s$ is timelike and one chooses the unit 4-velocity
$U_\mu:=\nabla_\mu s/|\nabla s|_g$, if the seam rule has no Hessian term
($\gamma\equiv 0$, so $g=\alpha h+\beta\,ds\otimes ds$), then the effective
stress--energy has vanishing anisotropic stress in the standard imperfect-fluid
decomposition relative to $U$.
\end{corollary}

\begin{proof}
When $\gamma\equiv 0$, the only generator contributions to $T^{\mathrm{eff}}$ are
the $\alpha$-term (proportional to $g$, hence isotropic) and the $\beta$-term
(proportional to $ds\otimes ds\propto U\otimes U$ when $\nabla s$ is timelike).
Both have zero trace-free spatial projection with respect to $U$, hence the
anisotropic stress $\pi_{\mu\nu}$ vanishes.
\end{proof}

\begin{corollary}[Irrotationality of the seam 4-velocity]
\label{cor:irrotational}
In any region where $\nabla s$ is timelike and one chooses the unit 4-velocity
$U_\mu:=\nabla_\mu s/|\nabla s|_g$, the associated flow is hypersurface-orthogonal
and has identically vanishing vorticity tensor.
\end{corollary}

\begin{proof}
Since $U$ is proportional to the exact 1-form $ds$, one has $U\wedge dU=0$.
By Frobenius' theorem, $U$ is hypersurface-orthogonal, which is equivalent to
vanishing vorticity.
\end{proof}

\smallskip
\noindent\textbf{Relation to disformal-type ans\"atze.}
If the Hessian generator is absent ($\gamma\equiv 0$), the seam rule reduces to a
conformal-plus-gradient form $g=\alpha\,h+\beta\,ds\otimes ds$, which is
structurally the same tensorial ansatz that appears in the disformal-metric
literature of scalar-tensor gravity; see, e.g.,~\cite{Bekenstein1993,CliftonFerreiraPadillaSkordis2012}.  The seam decomposition differs in two ways:
first, the coefficient functions are allowed to depend on the full list of
2-jet invariants~\eqref{eq:UD-invariants} (not only on $s$ and
$|\nabla s|^2$); and second, the additional Hessian generator provides an
extra direction in the space of natural order-$\le 2$ metric rules.

\smallskip
\noindent\textbf{Degrees of freedom and gravitational radiation (open question).}
Because the metric is generated from a scalar seam, seam-generated vacuum
metrics can be highly constrained: in the maximally symmetric families computed
here, the vacuum condition forces flatness.  Whether less symmetric seam
metrics (with genuinely nontrivial Hessian structure) can support nontrivial
vacuum Weyl curvature, or accommodate tensor gravitational radiation when viewed
as exact solutions rather than as perturbations around a seam-generated
background, remains to be investigated.

\subsection{Vacuum structure and the uniqueness of flat space}

Theorem~\ref{thm:vacuumon} shows that within the $O(n)$-invariant
Euclidean seam family, the unique vacuum solution is flat space.
This is a strong constraint: no nontrivial spherically symmetric
vacuum Riemannian 3-manifold exists in this family.  The Schwarzschild
geometry enters only as a vacuum spacetime (4D), where the time
direction provides the missing degree of freedom that allows a
non-trivial vacuum structure.

The conformal-seam vacuum condition (Theorem~\ref{thm:vacuum})
is similarly restrictive: $e^s$ must be harmonic, which by Liouville's
theorem forces $e^s$ to be constant on complete flat spaces.

Both results suggest a general principle: seam-generated metrics are
\emph{overdetermined} by their generators.  The generators encode so
much information that the vacuum condition eliminates almost all
non-trivial solutions, leaving only the geometrically simplest ones.
This is best read as evidence that the present scalar-generated ansatz classes
are highly restrictive.  That restriction is useful when the goal is to obtain
sharp no-go theorems or explicit effective-matter dictionaries, but it also
marks the boundary of what these families can represent.

\begin{remark}[Limitations of the real scalar seam]
\label{rem:complex-seam}
Three structural limitations of the real scalar seam are visible in the
results above:
\emph{(i)}~vacuum overdetermination (flatness rigidity in all symmetric
families),
\emph{(ii)}~a single functional degree of freedom versus the two
polarisation modes of gravitational radiation, and
\emph{(iii)}~identically vanishing vorticity
(Corollary~\ref{cor:irrotational}).
All three are traceable to the seam being a single real-valued function.
In Section~\ref{sec:complex} we show that promoting the seam to a
complex-valued field $\varphi=u+iv$ resolves each limitation:
limitation~(ii) via an independent phase degree of freedom,
limitation~(iii) via the cross-gradient $du\wedge dv$, and
limitation~(i) via the full six-generator span, which admits
genuinely curved Ricci-flat metrics
(Theorem~\ref{thm:complex-vacuum}).
\end{remark}

\subsection{Signature and the pre-geometric picture}

Theorem~\ref{thm:signature} shows that Lorentzian signature is not fixed
independently of the seam rule within the $O(n)$-invariant seam family.
Riemannian and Lorentzian
geometry are two regions of the same pointwise sign structure of $(A,P)$,
separated by the degeneration sphere where $P(r_0)=0$.  The
physical interpretation of this transition --- the radial direction
becoming timelike inside $r_0$ --- is not imposed but derived.

In the conformal family, the Weyl tensor vanishes identically; in the other two
symmetric families studied here, the vacuum condition forces flatness and hence
Weyl curvature does not arise in vacuum within these ansatz classes.

In this sense the seam scalar plays a role loosely analogous to Wheeler's
\emph{it from bit} idea~\cite{Wheeler1990}: metric signature is determined by
the seam rather than prescribed independently.

% ================================================================
\section{Energy Conditions as Seam Convexity (Conformal Family)}
\label{sec:energy}
% ================================================================

This section expands the ``effective seam matter'' viewpoint
\eqref{eq:Teff} by translating energy conditions into explicit inequalities on
the seam for the conformal family $g=e^{2s}\eta$ (cf.~\cite{HawkingEllis1973,Wald1984}).

\begin{proposition}[NEC as a seam convexity inequality (conformal family)]
\label{prop:nec}
For the conformal seam metric $g_{\mu\nu}=e^{2s}\eta_{\mu\nu}$ in $n=4$, and any
null vector $k$ (null with respect to $g$, equivalently to $\eta$), one has
\begin{equation}
  \label{eq:nec-ineq}
  T^{\mathrm{eff}}_{\mu\nu}k^\mu k^\nu \ge 0
  \quad\Longleftrightarrow\quad
  (k^\mu\pa_\mu s)^2 \ge k^\mu k^\nu\pa_\mu\pa_\nu s.
\end{equation}
\end{proposition}

\begin{proof}
Using~\eqref{eq:main} and $g_{\mu\nu}k^\mu k^\nu=0$ (conformal metrics preserve
the null cone), the $\alpha$-generator drops out and
\[
G_{\mu\nu}k^\mu k^\nu
=-2\,k^\mu k^\nu\pa_\mu\pa_\nu s + 2\,(k^\mu\pa_\mu s)^2.
\]
Dividing by $8\pi$ gives~\eqref{eq:nec-ineq}.
\end{proof}

\begin{lemma}[Reciprocal conformal factor identity]
\label{lem:reciprocal-hess}
Let $u:=e^{-s}$.  Then for any vector $k^\mu$,
\begin{equation}
  \label{eq:reciprocal-hess}
  k^\mu k^\nu\pa_\mu\pa_\nu u
  = u\bigl[(k^\mu\pa_\mu s)^2-k^\mu k^\nu\pa_\mu\pa_\nu s\bigr].
\end{equation}
\end{lemma}

\begin{proof}
$\pa_\mu u=-u\,\pa_\mu s$, hence
$\pa_\mu\pa_\nu u=u\,\pa_\mu s\,\pa_\nu s-u\,\pa_\mu\pa_\nu s$.
Contracting with $k^\mu k^\nu$ gives~\eqref{eq:reciprocal-hess}.
\end{proof}

\begin{corollary}[NEC as null-convexity of $u=e^{-s}$]
\label{cor:null-convex}
In the conformal family, the NEC inequality~\eqref{eq:nec-ineq} is equivalent to
\begin{equation}
  \label{eq:nullconvex}
  k^\mu k^\nu\pa_\mu\pa_\nu u \ge 0\quad\text{for all null }k.
\end{equation}
Equivalently, $u$ is convex when restricted to every affinely parametrised null
line in the flat background.
\end{corollary}

\begin{proof}
Lemma~\ref{lem:reciprocal-hess} gives
$k^\mu k^\nu\pa_\mu\pa_\nu u\ge 0\iff (k^\mu\pa_\mu s)^2\ge k^\mu k^\nu\pa_\mu\pa_\nu s$
since $u=e^{-s}>0$.
\end{proof}

\begin{lemma}[Reciprocal-Hessian decomposition]
\label{lem:recip-decomp}
For $g_{\mu\nu}=e^{2s}\eta_{\mu\nu}$ in $n=4$, set $u:=e^{-s}$ and
$H_{\mu\nu}:=\pa_\mu\pa_\nu u$.  Then
\begin{align}
  \label{eq:G-recip}
  G_{\mu\nu} &= \tfrac{2}{u}\,H_{\mu\nu}
    + \bigl(2\,\Box_\eta s+|\nn s|_\eta^2\bigr)\,\eta_{\mu\nu},\\
  \label{eq:R-recip}
  R_{\mu\nu} &= \tfrac{2}{u}\,H_{\mu\nu}
    - \bigl(\Box_\eta s+2\,|\nn s|_\eta^2\bigr)\,\eta_{\mu\nu}.
\end{align}
In particular, for null~$k$ both decompositions give the same contraction:
$G_{\mu\nu}k^\mu k^\nu=R_{\mu\nu}k^\mu k^\nu=\frac{2}{u}\,H_{\mu\nu}k^\mu k^\nu$.
\end{lemma}

\begin{proof}
In~\eqref{eq:main}, the $\gamma$- and $\beta$-generators combine via
Lemma~\ref{lem:reciprocal-hess}:
$-2\,\pa_\mu\pa_\nu s+2\,(\pa_\mu s)(\pa_\nu s)=\frac{2}{u}\,H_{\mu\nu}$.
The $\alpha$-generator satisfies $e^{-2s}g_{\mu\nu}=\eta_{\mu\nu}$,
giving~\eqref{eq:G-recip}.  For~\eqref{eq:R-recip}, substitute
Proposition~\ref{prop:ricci} with $n=4$ and apply the same identity.
For null~$k$, $\eta_{\mu\nu}k^\mu k^\nu=0$ kills both flat-metric terms.
\end{proof}

\begin{corollary}[Constant curvature from an affine reciprocal seam]
\label{cor:affine-u-cc}
Let $g_{\mu\nu}=e^{2s}\eta_{\mu\nu}$ and set $u:=e^{-s}$.  If on a domain
$\Omega\subset\mathbb{R}^{1,3}$ with $u>0$ one has an affine reciprocal factor
$u(x)=k_\mu x^\mu+c$ for a constant covector $k$ and constant $c$, then
\begin{equation}
  \label{eq:affine-u-einstein}
  G_{\mu\nu}=3\,|k|_\eta^2\,g_{\mu\nu}\qquad\text{on }\Omega.
\end{equation}
Equivalently, $g$ is an Einstein metric with cosmological constant
$\Lambda=-3|k|_\eta^2$ on $\Omega$ (de~Sitter type if $k$ is timelike and
anti--de~Sitter type if $k$ is spacelike).  Moreover, the NEC
inequality~\eqref{eq:nec-ineq} is saturated for all null vectors.
\end{corollary}

\begin{proof}
If $u$ is affine then $\pa_\mu u=k_\mu$ and $\pa_\mu\pa_\nu u=0$.  Since
$s=-\log u$, one has $\pa_\mu s=-k_\mu/u$ and $\pa_\mu\pa_\nu s=k_\mu k_\nu/u^2$,
so the $\gamma$- and $\beta$-generators in~\eqref{eq:main} cancel:
$-2\,\pa_\mu\pa_\nu s+2(\pa_\mu s)(\pa_\nu s)=0$.
Also $\Box_\eta s=|k|_\eta^2/u^2$ and $|\nabla s|_\eta^2=|k|_\eta^2/u^2$.
Since $e^{-2s}=u^2$, the $\alpha$-generator becomes
$u^2\,(3|k|_\eta^2/u^2)g_{\mu\nu}=3|k|_\eta^2 g_{\mu\nu}$.
Finally, by Lemma~\ref{lem:reciprocal-hess} one has
$(k^\mu\pa_\mu s)^2-k^\mu k^\nu\pa_\mu\pa_\nu s=(k^\mu k^\nu\pa_\mu\pa_\nu u)/u$;
for affine $u$ the Hessian vanishes, hence~\eqref{eq:nec-ineq} holds with equality.
\end{proof}

\begin{proposition}[Averaged NEC inequality (necessary condition)]
\label{prop:anec}
If $u=e^{-s}$ is null-convex (equivalently, if the NEC holds), then
\begin{equation}
  \label{eq:anec}
  3\,\pa^2_t u + \Delta_x u \ge 0,
\end{equation}
where $\Delta_x$ is the spatial Laplacian in the flat background.
The converse fails in general: \eqref{eq:anec} is necessary but not sufficient
for the NEC.
\end{proposition}

\begin{proof}
For each unit $\mathbf{v}\in S^2$, the vector $k_{\mathbf{v}}=(1,\mathbf{v})$ is
null.  Null-convexity gives
$k_{\mathbf{v}}^\mu k_{\mathbf{v}}^\nu\pa_\mu\pa_\nu u\ge 0$, i.e.
$\pa^2_t u + 2\mathbf{v}\cdot\nabla\pa_t u + \mathbf{v}^i\mathbf{v}^j\pa_i\pa_j u\ge 0$.
Averaging over $\mathbf{v}\in S^2$ and using
$\langle v^i\rangle=0$, $\langle v^iv^j\rangle=\tfrac{1}{3}\delta^{ij}$ yields
$\pa^2_t u+\tfrac{1}{3}\Delta_x u\ge 0$, hence~\eqref{eq:anec}.
Non-sufficiency follows by a pointwise counterexample at the level of Hessians.
\end{proof}

\begin{theorem}[Rigidity under a dynamical assumption]
\label{thm:rigid}
Let $g_{\mu\nu}=e^{2s}\eta_{\mu\nu}$ and $u=e^{-s}$.  Suppose:
\begin{enumerate}[label=(\roman*),noitemsep]
\item the NEC holds (equivalently, $u$ is null-convex);
\item \emph{(additional dynamical assumption)} $u$ satisfies the wave equation
  $\Box_\eta u=0$;
\item $u$ is bounded on all of $\mathbb{R}^{1,3}$.
\end{enumerate}
Then $u=\mathrm{const}$ and $g=\mathrm{const}\cdot\eta$.
\end{theorem}

\begin{proof}
Proposition~\ref{prop:anec} gives $\pa_t^2 u+\tfrac{1}{3}\Delta_x u\ge 0$.
From $\Box_\eta u=0$ one has $\Delta_x u=\pa_t^2 u$.  Hence
$\tfrac{4}{3}\pa_t^2 u\ge 0$, so $\pa_t^2 u\ge 0$ and thus $\Delta_x u\ge 0$.
For each fixed $t$, $u(\cdot,t)$ is a bounded subharmonic function on
$\mathbb{R}^3$, hence constant.  The wave equation then forces $u$ to be affine
in $t$, and boundedness forces it to be constant.
\end{proof}

\noindent The assumption $\Box_\eta u=0$ is the free massless wave equation on
the flat background.  It is logically independent of the conformal-seam vacuum
condition $\Box_\eta(e^{s})=0$ from Theorem~\ref{thm:vacuum}.  Expanding
$\Box_\eta(e^{-s})=0$ yields the equivalent scalar relation
$\Box_\eta s = |\nabla s|_\eta^2$.

\smallskip
\noindent\textbf{Compatibility note.}
While neither wave equation implies the other, requiring both
$\Box_\eta(e^{s})=0$ and $\Box_\eta(e^{-s})=0$ is highly restrictive: it forces
the seam to have null gradient, hence places the conformal factor on a
null-propagating constraint surface (Proposition~\ref{prop:dual-wave}).

\begin{proposition}[Dual wave equations force a null seam]
\label{prop:dual-wave}
Let $s\in C^\infty(\mathbb{R}^{1,3})$.  If both
\begin{equation}
  \label{eq:dual-wave-assumptions}
  \Box_\eta(e^{s})=0
  \quad\text{and}\quad
  \Box_\eta(e^{-s})=0,
\end{equation}
then $|\nabla s|_\eta^2\equiv 0$ and $\Box_\eta s\equiv 0$.
\end{proposition}

\begin{proof}
Expanding $\Box_\eta(e^{s})=e^{s}(\Box_\eta s+|\nabla s|_\eta^2)$ gives
$\Box_\eta s=-|\nabla s|_\eta^2$.  Also
$\Box_\eta(e^{-s})=e^{-s}(|\nabla s|_\eta^2-\Box_\eta s)$.
Substituting yields $\Box_\eta(e^{-s})=2e^{-s}|\nabla s|_\eta^2$; hence
$|\nabla s|_\eta^2=0$, and then $\Box_\eta s=0$.
\end{proof}

\begin{remark}
One large class of examples satisfying~\eqref{eq:dual-wave-assumptions} is given
by seams of the form $s=f(\phi)$ where $\phi$ is a null coordinate on Minkowski
space (so $|\nabla \phi|_\eta^2=0$), for instance $\phi=k_\mu x^\mu$ with a
constant null covector $k$.  Then both $e^{\pm s}$ depend on a single null
variable and automatically solve the flat wave equation.
\end{remark}

\begin{proposition}[Energy conditions via the reciprocal Hessian]
\label{prop:ec-dictionary}
For $g_{\mu\nu}=e^{2s}\eta_{\mu\nu}$ in $n=4$, let $u:=e^{-s}$,
$H_{\mu\nu}:=\pa_\mu\pa_\nu u$, and define the scalars
\[
  P:=2\,\Box_\eta s+|\nn s|^2_\eta,
  \qquad
  Q:=\Box_\eta s+2\,|\nn s|^2_\eta.
\]
The pointwise energy conditions, stated for every $\eta$-unit timelike vector
$n^\mu$ ($\eta_{\mu\nu}n^\mu n^\nu=-1$), are:
\begin{enumerate}[label=\textup{(\roman*)},noitemsep]
\item \textup{NEC} $\;\iff\; H_{\mu\nu}k^\mu k^\nu\ge 0$ for all null~$k$
  \;\textup{(null-convexity of $u$)};
\item \textup{WEC} $\;\iff$ NEC and\;
  $\dfrac{2}{u}\,H_{\mu\nu}n^\mu n^\nu \ge P$ for all such~$n$;
\item \textup{SEC} $\;\iff\;
  \dfrac{2}{u}\,H_{\mu\nu}n^\mu n^\nu \ge -Q$ for all such~$n$
  \;\textup{(implies NEC)};
\item \textup{DEC} $\;\iff$ WEC and\;
  $\dfrac{2}{u}\,H_{\mu\nu}k^\mu \ell^\nu
    +P\,\eta_{\mu\nu}k^\mu\ell^\nu\ge 0$
  for all future-directed null~$k,\ell$.
\end{enumerate}
\end{proposition}

\begin{proof}
By Lemma~\ref{lem:recip-decomp},
$G_{\mu\nu}=\frac{2}{u}H_{\mu\nu}+P\eta_{\mu\nu}$ and
$R_{\mu\nu}=\frac{2}{u}H_{\mu\nu}-Q\eta_{\mu\nu}$.

(i) For null~$k$, $G_{\mu\nu}k^\mu k^\nu=\frac{2}{u}H_{\mu\nu}k^\mu k^\nu$
since $\eta_{\mu\nu}k^\mu k^\nu=0$.  With $u>0$, NEC is null-convexity of~$u$.

(ii) WEC requires $G_{\mu\nu}v^\mu v^\nu\ge 0$ for all $g$-timelike~$v$.
A $g$-unit timelike vector satisfies $g_{\mu\nu}v^\mu v^\nu
=e^{2s}\eta_{\mu\nu}v^\mu v^\nu=-1$, so $v=u\,n$ with
$\eta(n,n)=-1$ parametrises all such vectors.  Hence
$G_{\mu\nu}v^\mu v^\nu
  =u^2\bigl[\tfrac{2}{u}H(n,n)-P\bigr]$,
and since $u^2>0$, WEC reduces to
$\tfrac{2}{u}H(n,n)\ge P$ for every $\eta$-unit timelike~$n$.
The left-hand side is minimised when $n$ is aligned with the
$\eta$-timelike eigenvector of~$H_{\mu\nu}$
(the unique direction that minimises $H(n,n)$ over the unit
hyperboloid), so the binding constraint is the stated inequality
evaluated at that eigenvector.

(iii) SEC requires $R_{\mu\nu}v^\mu v^\nu\ge 0$ for all timelike~$v$.
With $v=u\,n$:
$R_{\mu\nu}v^\mu v^\nu=u^2[\tfrac{2}{u}H(n,n)+Q]$, giving~(iii).
Null vectors are limits of unit timelike vectors, and
$H(k,k)\ge 0$ follows from the limit, so SEC implies NEC.

(iv) DEC is equivalent to $G_{\mu\nu}v^\mu w^\nu\ge 0$ for all future-directed
causal $v,w$~\cite[Prop.~4.3.1]{HawkingEllis1973}.  Every future-directed
causal vector is a non-negative linear combination of future-directed null
vectors.  Hence $G_{\mu\nu}v^\mu w^\nu\ge 0$ for all causal pairs iff it holds
for all null pairs $k,\ell$.  The timelike--timelike case (WEC) then follows by
restriction, and the null--null condition is as stated.
\end{proof}

\begin{corollary}[Energy condition dictionary for temporal seams]
\label{cor:ec-temporal}
For a purely temporal seam $s=g(t)$, set $u:=e^{-g(t)}$.
Primes denote derivatives with respect to~$t$.  Then:
\begin{align}
  \textup{NEC}\;(=\textup{WEC})\;
    &\iff g''\le (g')^2
    \iff u''\ge 0,
    \label{eq:nec-wec-temporal}\\
  \textup{SEC}\;
    &\iff g''\le 0
    \iff u''\ge (u')^2/u,
    \label{eq:sec-temporal}\\
  \textup{DEC}\;
    &\iff {-2(g')^2}\le g''\le (g')^2
    \iff 0\le u''\le 3(u')^2/u.
    \label{eq:dec-temporal}
\end{align}
NEC and WEC coincide because the energy density
$\rho=3(g')^2 e^{-2g}/(8\pi)$ is automatically non-negative.
SEC requires seam concavity ($g''\le 0$) and is strictly stronger than NEC.
DEC adds a lower bound on~$g''$, constraining the deceleration rate.
\end{corollary}

\begin{proof}
For $s=g(t)$: $H_{00}=u''=(g'^2-g'')u$, all other $H_{\mu\nu}=0$;
$P=-2g''-(g')^2$; $Q=-g''-2(g')^2$; $|\nn u|^2_\eta=-(g')^2u^2$.

(i--ii) NEC: $H(k,k)=(k^0)^2 u''\ge 0$ for null $k=(1,\hat{v})$, giving
$u''\ge 0$, i.e.\ $g''\le(g')^2$.
WEC adds $\frac{2}{u}H(n,n)\ge P$ for $\eta$-unit timelike~$n$.
At the rest frame $n=(1,0,0,0)$: $\frac{2}{u}u''=2(g'^2-g'')$ and
$P=-2g''-(g')^2$; the bound reads $2(g'^2-g'')\ge -2g''-(g')^2$, i.e.\
$3(g')^2\ge 0$---always satisfied.  Boosted observers give a weaker constraint
(controlled by NEC), so WEC $\iff$ NEC.

(iii) SEC adds $\frac{2}{u}H(n,n)\ge -Q$.  At the rest frame:
$2(g'^2-g'')\ge g''+2(g')^2$, i.e.\ $-3g''\ge 0$.  Hence SEC $\iff g''\le 0$.
In $u$-variables: $g''\le 0\iff u''\ge (u')^2/u$.

(iv) DEC adds the null--null condition for $k=(1,\hat{m})$,
$\ell=(1,\hat{n})$ future-null:
$\frac{2u''}{u}+P(-1+\hat{m}\cdot\hat{n})\ge 0$.
The most constraining case is $\hat{m}\cdot\hat{n}=-1$ (antiparallel), giving
$2(g'^2-g'')+2(2g''+(g')^2)\ge 0$, i.e.\
$4(g')^2+2g''\ge 0$, so $g''\ge -2(g')^2$.
Combined with NEC: $-2(g')^2\le g''\le (g')^2$.
\end{proof}

\begin{corollary}[Past vanishing scale factor under NEC for temporal seams]
\label{cor:temporal-past}
Let $s=g(t)$ be a temporal seam defined on $(-\infty,t_0]$ and set
$a(t):=e^{g(t)}$ and $u(t):=e^{-g(t)}=1/a(t)$.  If the NEC holds and the seam is
expanding at some time $t_0$ (i.e. $g'(t_0)>0$, equivalently $u'(t_0)<0$), then
\begin{equation}
  \label{eq:temporal-past-limit}
  \lim_{t\to-\infty}u(t)=+\infty
  \qquad\text{and hence}\qquad
  \lim_{t\to-\infty}a(t)=0.
\end{equation}
In particular, $g(t)\to-\infty$ as $t\to-\infty$.
\end{corollary}

\begin{proof}
By Corollary~\ref{cor:ec-temporal}, NEC is equivalent to $u''\ge 0$, so $u$
is convex.  Convexity implies that for all $t\le t_0$,
$u(t)\ge u(t_0)+u'(t_0)(t-t_0)$.  Since $u'(t_0)<0$, the right-hand side tends
to $+\infty$ as $t\to-\infty$, hence $u(t)\to+\infty$ and $a(t)=1/u(t)\to 0$.
\end{proof}

\begin{remark}
Corollary~\ref{cor:temporal-past} is a one-dimensional convexity consequence in
the seam time coordinate $t$.  By itself it does not assert curvature blow-up
or geodesic incompleteness; such conclusions require additional hypotheses on
the spacetime interpretation of $t$ and on the behaviour of curvature
invariants.
\end{remark}

% ================================================================
\section{Result 4: Wavefunction Seam and Klein--Gordon Compatibility}
\label{sec:wavefunction}
% ================================================================

The conformal seam framework admits a natural quantum-mechanical
interpretation when the conformal factor is identified with the modulus
of a complex scalar field.

Let $\psi=|\psi|\,e^{i\theta}$ be a complex scalar field on Minkowski background
$(\mathbb{R}^{1,3},\eta)$, with $|\psi|>0$ and $\theta$ real-valued.
Define the \emph{wavefunction seam} by
\begin{equation}
  \label{eq:wf-seam}
  s:=\log|\psi|,
  \qquad\text{so that}\qquad
  g_{\mu\nu}=|\psi|^2\,\eta_{\mu\nu}.
\end{equation}
This is the conformal family of Result~1 with $e^{2s}=|\psi|^2$.
The reciprocal factor of Section~\ref{sec:energy} becomes $u=1/|\psi|$.

\begin{proposition}[Seam dictionary in wavefunction variables]
\label{prop:wf-dictionary}
For the wavefunction seam~\eqref{eq:wf-seam}:
\begin{enumerate}[label=\textup{(\roman*)},noitemsep]
\item \textup{Vacuum} \textup{(Theorem~\ref{thm:vacuum})}: $G_{\mu\nu}=0$
  if and only if $|\psi|$ is harmonic in the flat
  background: $\Box_\eta|\psi|=0$.
\item \textup{NEC} \textup{(Corollary~\ref{cor:null-convex})}: the null
  energy condition holds if and only if $1/|\psi|$ is convex along every
  null line:
  \begin{equation}
    \label{eq:wf-nec}
    k^\mu k^\nu\,\pa_\mu\pa_\nu(1/|\psi|)\ge 0
    \quad\text{for all null }k.
  \end{equation}
\item \textup{Harmonic ladder} \textup{(Corollary~\ref{cor:harmonic-hierarchy})}:
  the scalar curvature vanishes if and only if $|\psi|$ is harmonic;
  the isotropic generator vanishes if and only if $|\psi|^{1/2}$ is harmonic.
\end{enumerate}
\end{proposition}

\begin{proof}
Direct substitution of $s=\log|\psi|$ into the cited results.
\end{proof}

The vacuum condition acquires a sharp physical meaning when $\psi$
satisfies a field equation.

\begin{theorem}[Klein--Gordon compatibility]
\label{thm:kg-compat}
Let $\psi=|\psi|\,e^{i\theta}$ satisfy the Klein--Gordon equation
$(\Box_\eta+m^2)\psi=0$ on a region where $|\psi|>0$.
Then the seam-geometric vacuum condition $G_{\mu\nu}=0$
for the wavefunction metric $g=|\psi|^2\eta$ holds if and only if
the phase gradient satisfies the eikonal equation:
\begin{equation}
  \label{eq:eikonal}
  |\nn\theta|^2_\eta = m^2.
\end{equation}
\end{theorem}

\begin{proof}
Substituting $\psi=|\psi|\,e^{i\theta}$ into the Klein--Gordon equation
and separating real parts gives
\begin{equation}
  \label{eq:kg-real}
  \Box_\eta|\psi| = |\psi|\,(|\nn\theta|^2_\eta-m^2).
\end{equation}
(The imaginary part gives the continuity equation
$\pa_\mu(|\psi|^2\,\nn^\mu\theta)=0$.)
By Proposition~\ref{prop:wf-dictionary}(i), the vacuum condition is
$\Box_\eta|\psi|=0$.  Since $|\psi|>0$,
\eqref{eq:kg-real} then gives~\eqref{eq:eikonal}.
Conversely, \eqref{eq:eikonal} substituted into~\eqref{eq:kg-real}
yields $\Box_\eta|\psi|=0$.
\end{proof}

\begin{remark}[Eikonal regime and the classical limit]
\label{rem:eikonal}
The eikonal equation $|\nn\theta|^2_\eta=m^2$ is precisely the condition
defining the WKB (geometric-optics) limit of the Klein--Gordon field.
In this regime the phase surfaces propagate as classical wavefronts with
4-momentum $p_\mu=\pa_\mu\theta$ satisfying the mass-shell relation
$p^\mu p_\mu=-m^2$.
For the massless case $m=0$, the eikonal condition forces
$\nn\theta$ to be null, corresponding to light-like propagation.
Theorem~\ref{thm:kg-compat} therefore identifies the seam-geometric vacuum
as the classical-propagation limit of the underlying quantum field:
the wavefunction metric is vacuum-Einstein precisely when the
field enters the WKB regime.
\end{remark}

\begin{remark}[Quantum potential]
\label{rem:quantum-potential}
In the Bohmian/Madelung formulation of relativistic quantum mechanics,
the quantity $\Box_\eta|\psi|/|\psi|$ plays the role of a quantum
potential~\cite{ShojaiShojai1998}.
From~\eqref{eq:kg-real}, the Klein--Gordon real part is the
relativistic quantum Hamilton--Jacobi equation
\[
  |\nn\theta|^2_\eta = m^2 + \frac{\Box_\eta|\psi|}{|\psi|}.
\]
The seam-geometric vacuum condition $\Box_\eta|\psi|=0$ is therefore
equivalent to the vanishing of the quantum potential, and the eikonal
equation~\eqref{eq:eikonal} emerges as the zero--quantum-potential limit.
\end{remark}

\begin{remark}[Relation to prior work]
\label{rem:bohmian-prior}
The idea that $|\psi|$ (or a quantum potential derived from it) sources or
modifies spacetime geometry has appeared in the Bohmian gravity
literature---notably in the Shojai--Shojai programme on geometrisation of
Bohmian mechanics~\cite{ShojaiShojai1998} and in Chowdhury et al., who couple
$|\Phi|$ to the Ricci scalar via disformal
transformations~\cite{Chowdhury2021}.
The contribution here is the structural
observation that (i)~the Einstein tensor of the wavefunction metric
decomposes into the three seam generators with
$\psi$-dependent coefficients, (ii)~the vacuum condition has a sharp
characterisation as the eikonal equation for the phase
(equivalently, the vanishing of the quantum potential), and (iii)~the
energy conditions reduce to null-convexity of $1/|\psi|$.
\end{remark}

\subsection{Physical corollaries of the wavefunction seam}
\label{subsec:wf-physics}

Combining the results of Sections~\ref{sec:conformal},
\ref{sec:energy}, and the present section yields several physical
observations that are immediate consequences of the established
formalism.

\begin{corollary}[Scalar curvature as quantum potential]
\label{cor:R-quantum}
For the wavefunction metric $g_{\mu\nu}=|\psi|^2\,\eta_{\mu\nu}$ in $n=4$,
the scalar curvature satisfies
\begin{equation}
  \label{eq:R-Q}
  R = -\frac{6}{|\psi|^2}\,Q,
  \qquad
  Q := \frac{\Box_\eta|\psi|}{|\psi|},
\end{equation}
where $Q$ is the relativistic quantum potential of
Remark~\ref{rem:quantum-potential}.
In particular, $R=0$ if and only if the quantum potential vanishes,
and signed deviations of $R$ correspond directly to signed deviations
of $Q$ from the classical (WKB) limit.
\end{corollary}

\begin{proof}
From Corollary~\ref{cor:harmonic-hierarchy}(i), $R=0$ iff
$\Box_\eta(e^s)=0$, i.e.\ $\Box_\eta|\psi|=0$.
More explicitly, Proposition~\ref{prop:ricciscalar} with $s=\log|\psi|$ gives
$R=-6e^{-2s}(\Box_\eta s+|\nn s|^2_\eta)=-6|\psi|^{-2}\cdot
\Box_\eta|\psi|/|\psi|$,
using $\Box_\eta(e^s)=e^s(\Box_\eta s+|\nn s|^2)$.
\end{proof}

\begin{corollary}[Metric degeneration at wavefunction nodes]
\label{cor:nodes}
Let $\psi=|\psi|\,e^{i\theta}$ with $|\psi|\ge 0$.  At any node
$x_0$ where $|\psi(x_0)|=0$, the wavefunction metric
$g_{\mu\nu}=|\psi|^2\,\eta_{\mu\nu}$ degenerates: all distances measured
by $g$ shrink to zero and the reciprocal factor $u=1/|\psi|$ diverges.
In the temporal reduction $|\psi|=a(t)$, this coincides with the
$a\to 0$ singularity of Corollary~\ref{cor:temporal-past}.  Every
zero of the amplitude thus corresponds to a conformal singularity of
the effective spacetime.
\end{corollary}

\begin{proof}
$g_{\mu\nu}(x_0)=|\psi(x_0)|^2\eta_{\mu\nu}=0$.  In the temporal case,
$a(t)=|\psi(t)|$, and the FLRW-type metric $g=-\mathrm{d}t^2\cdot a^2+\cdots$
degenerates at $a=0$.
\end{proof}

\begin{corollary}[No event horizons in the wavefunction metric]
\label{cor:no-horizons}
The wavefunction metric $g=|\psi|^2\eta$ cannot admit an event horizon.
\end{corollary}

\begin{proof}
By Corollary~\ref{cor:no-lensing}, the unparameterised null geodesics of $g$
are the straight null lines of $\eta$.  An event horizon requires the existence
of a trapped region from which no future-directed null geodesic can escape to
infinity.  Since every null ray of $g$ traces a straight line in the background,
no such trapping is possible.
\end{proof}

\begin{remark}[NEC violation near classical turning points]
\label{rem:turning-points}
By Proposition~\ref{prop:wf-dictionary}(ii), the NEC holds iff $1/|\psi|$ is
convex along every null line.  For a Klein--Gordon field tunnelling through a
classically forbidden region, the WKB amplitude decays as
$|\psi|\sim e^{-\kappa x}$, giving $1/|\psi|\sim e^{\kappa x}$, which is
convex; the NEC is satisfied inside the barrier.  At classical turning points,
however, the amplitude changes character from oscillatory to evanescent and may
develop inflection points in $1/|\psi|$, violating the NEC locally.  Since NEC
violation is the necessary ingredient for exotic phenomena such as traversable
wormhole throats in classical GR~\cite{MorrisThorne1988}, the transition
region of a tunnelling wavefunction carries a direct geometric signature within
this framework.
\end{remark}

\begin{remark}[Quantum interference breaks the vacuum condition]
\label{rem:interference}
The vacuum condition $\Box_\eta|\psi|=0$ is nonlinear in $\psi$ through
the modulus map.  If $\psi_1$ and $\psi_2$ each satisfy the vacuum condition
($\Box_\eta|\psi_1|=\Box_\eta|\psi_2|=0$), their quantum superposition
$\psi=\psi_1+\psi_2$ generically does not: $|\psi_1+\psi_2|$ involves the
interference term $\mathrm{Re}(\psi_1\bar\psi_2)$, and the resulting
amplitude is not harmonic.  By Corollary~\ref{cor:R-quantum}, the non-vanishing
$\Box_\eta|\psi|/|\psi|$ produces scalar curvature $R\neq 0$ localised at
the interference fringes.  In this sense, quantum superposition of
individually vacuum (WKB-limit) states generates effective stress-energy
and curvature in the wavefunction metric.
\end{remark}

\begin{corollary}[Bound states are never vacuum]
\label{cor:bound-nonvacuum}
Let $\psi(t,x)=\phi(x)\,e^{-iEt}$ be a stationary bound state on
$\mathbb{R}^{1,3}$ with $|\phi(x)|\to 0$ as $|x|\to\infty$ and
$\phi\not\equiv 0$.  Then the wavefunction metric
$g_{\mu\nu}=|\phi|^2\,\eta_{\mu\nu}$ is not vacuum:
$G_{\mu\nu}\neq 0$.
\end{corollary}

\begin{proof}
Since $|\psi|=|\phi(x)|$ is time-independent, the vacuum condition
$\Box_\eta|\psi|=0$ reduces to $\Delta|\phi|=0$ on $\mathbb{R}^3$.
But $|\phi|\ge 0$ with $|\phi|\to 0$ at infinity.  By the maximum
principle for harmonic functions, a non-negative harmonic function that
vanishes at infinity on $\mathbb{R}^3$ must be identically zero.  Since
$\phi\not\equiv 0$, the vacuum condition fails.
\end{proof}

\begin{remark}[Physical reading of Corollary~\ref{cor:bound-nonvacuum}]
\label{rem:bound-reading}
Corollary~\ref{cor:bound-nonvacuum} shows that spatial quantum confinement
is incompatible with a vacuum wavefunction metric.
Only completely delocalised states (such as plane waves with
$|\psi|=\mathrm{const}$) can yield $G_{\mu\nu}=0$.
For any localised bound state, the quantum potential is necessarily non-zero
and the effective spacetime carries intrinsic curvature.
\end{remark}

\subsection{Self-consistent wavefunction seam}
\label{subsec:self-consistent}

So far the wavefunction $\psi$ satisfies the Klein--Gordon equation on the
\emph{flat} background $\eta$.  The metric $g=|\psi|^2\eta$ is then read off
as a derived object; the field does not propagate on the geometry it generates.
A natural strengthening is to require \emph{self-consistency}: the field
satisfies Klein--Gordon on $g$ itself,
\begin{equation}
  \label{eq:self-consistent}
  (\Box_g+m^2)\psi=0,
  \qquad
  g_{\mu\nu}=|\psi|^2\,\eta_{\mu\nu}.
\end{equation}

\begin{proposition}[Conformal wave operator]
\label{prop:conformal-box}
For $g_{\mu\nu}=e^{2s}\eta_{\mu\nu}$ in $n=4$, the minimally coupled
d'Alembertian satisfies
\begin{equation}
  \label{eq:conformal-box}
  \Box_g\phi
  = e^{-2s}\bigl[\Box_\eta\phi+2\,(\pa^\mu_\eta s)\,\pa_\mu\phi\bigr]
\end{equation}
for any scalar $\phi$.
\end{proposition}

\begin{proof}
$g^{\mu\nu}=e^{-2s}\eta^{\mu\nu}$ and $\sqrt{|g|}=e^{4s}\sqrt{|\eta|}$.
Expanding the covariant divergence
$\Box_g\phi=|g|^{-1/2}\pa_\mu(\sqrt{|g|}\,g^{\mu\nu}\pa_\nu\phi)$
gives~\eqref{eq:conformal-box}.
\end{proof}

With $s=\log|\psi|$, Proposition~\ref{prop:conformal-box} converts the
self-consistent equation~\eqref{eq:self-consistent} into a nonlinear wave
equation on the flat background:
\begin{equation}
  \label{eq:sc-flat}
  \Box_\eta\psi
  + \frac{2\,\pa^\mu|\psi|}{|\psi|}\,\pa_\mu\psi
  + m^2|\psi|^2\,\psi = 0.
\end{equation}

\begin{proposition}[Madelung form of the self-consistent system]
\label{prop:sc-madelung}
Under the polar decomposition $\psi=A\,e^{i\theta}$ with
$A=|\psi|>0$, equation~\eqref{eq:sc-flat} splits into:
\begin{enumerate}[label=\textup{(\roman*)},noitemsep]
\item \textbf{Modified Hamilton--Jacobi (real part):}
\begin{equation}
  \label{eq:sc-real}
  Q + 2\,|\nn\!\log A|_\eta^2
  = |\nn\theta|_\eta^2 - m^2 A^2,
\end{equation}
where $Q=\Box_\eta A/A$ is the quantum potential.
\item \textbf{Modified continuity (imaginary part):}
\begin{equation}
  \label{eq:sc-continuity}
  \pa_\mu\!\bigl(A^4\,\pa^\mu_\eta\theta\bigr)=0.
\end{equation}
\end{enumerate}
\end{proposition}

\begin{proof}
Substitute $\psi=Ae^{i\theta}$ into~\eqref{eq:sc-flat} and separate
real and imaginary parts.
The real part yields~\eqref{eq:sc-real} after dividing by~$A$.
For the imaginary part, one obtains
$4\,(\pa^\mu A)(\pa_\mu\theta)+A\,\Box_\eta\theta=0$, which is
$A^{-3}\,\pa_\mu(A^4\,\pa^\mu\theta)=0$.
\end{proof}

\begin{remark}[Comparison with the flat-background system]
\label{rem:sc-comparison}
On flat background (Section~\ref{sec:wavefunction}), the Klein--Gordon
Madelung system reads $Q=|\nn\theta|^2-m^2$ and
$\pa_\mu(A^2\pa^\mu\theta)=0$.  The self-consistent
version~\eqref{eq:sc-real}--\eqref{eq:sc-continuity} introduces two
modifications:
\begin{itemize}[noitemsep]
\item the mass term $m^2$ is replaced by the amplitude-dependent
  $m^2A^2$, producing a \emph{nonlinear dispersion relation};
\item the conserved current is $A^4\pa^\mu\theta$ instead of
  $A^2\pa^\mu\theta$; equivalently, the continuity equation is
  $\pa_\mu(\sqrt{|g|}\,g^{\mu\nu}\pa_\nu\theta)=0$, i.e.\ the current
  is the natural $g$-divergence-free vector.
\end{itemize}
\end{remark}

\begin{remark}[Comparison with the Schr\"odinger--Newton system]
\label{rem:sn-comparison}
The standard model for a self-gravitating nonrelativistic quantum field
is the Schr\"odinger--Newton (SN) equation: $i\hbar\pa_t\psi=-(\hbar^2/2m)\Delta\psi+m\Phi_N\psi$
with $\Delta\Phi_N=4\pi G m|\psi|^2$, where the Newtonian potential
$\Phi_N$ is sourced by the probability density.  The self-consistent
seam system~\eqref{eq:self-consistent} differs from SN in three structural
respects:
\begin{enumerate}[label=\textup{(\roman*)},noitemsep]
\item \emph{Fully relativistic:} the field satisfies the Klein--Gordon
  (not Schr\"odinger) equation on the curved metric, so no
  nonrelativistic limit is assumed.
\item \emph{Conformal coupling:} the gravitational backreaction
  enters conformally ($g=|\psi|^2\eta$), not through a Newtonian
  potential; this replaces the Poisson equation with a nonlinear wave
  equation~\eqref{eq:sc-flat}.
\item \emph{Current:} the conserved current is $A^4\pa^\mu\theta$
  (the $g$-natural current) versus $A^2\nn\theta$ in SN; the extra
  $A^2$ factor encodes the volume element of the self-generated
  geometry.
\end{enumerate}
In the nonrelativistic, weak-field limit ($A\approx 1+\epsilon$ with
$|\epsilon|\ll 1$, $|\pa_t|\ll m$), the self-consistent system
reduces to a Schr\"odinger equation with a nonlinear potential
$\propto|\psi|^2$, which has the same scaling as SN but with a
conformal coupling constant in place of $Gm^2$.  A detailed
comparison of the two nonlinear eigenvalue problems (Emden--Fowler
versus Choquard--Pekar) is deferred to future work.
\end{remark}

\begin{corollary}[Self-consistent vacuum fixed point]
\label{cor:sc-vacuum}
A solution of the self-consistent system~\eqref{eq:self-consistent} with
$|\psi|>0$ has vanishing Einstein tensor ($G_{\mu\nu}=0$, equivalently
$Q=0$) if and only if
\begin{equation}
  \label{eq:sc-eikonal}
  |\nn\theta|^2_\eta = m^2\,|\psi|^2 + 2\,|\nn\!\log|\psi||^2_\eta.
\end{equation}
For a constant-amplitude solution $|\psi|=A_0=\mathrm{const}$, this
reduces to $|\nn\theta|^2_\eta=m^2A_0^2$: the standard mass-shell
relation with effective mass $m_{\mathrm{eff}}=mA_0$, consistent with
the rescaled metric $g=A_0^2\eta$.
\end{corollary}

\begin{proof}
Set $Q=0$ in~\eqref{eq:sc-real}.  For $A=A_0=\mathrm{const}$,
$|\nn\log A|=0$ and the wavefunction metric is $g=A_0^2\eta$, on which
$(\Box_g+m^2)\psi=0$ becomes $\Box_\eta\psi+m^2A_0^2\psi=0$.
\end{proof}

\begin{remark}[Physical interpretation]
\label{rem:sc-physics}
The self-consistent system~\eqref{eq:self-consistent} closes the loop
between the wavefunction and its generated geometry: the metric is
sourced by $|\psi|$, and $\psi$ propagates on that metric.  Fixed-point
solutions---configurations satisfying $G_{\mu\nu}=0$ together
with~\eqref{eq:sc-flat}---are simultaneously vacuum geometries \emph{and}
WKB-limit fields, but with the modified eikonal
relation~\eqref{eq:sc-eikonal} replacing the standard mass-shell.  The
amplitude dependence in $m^2|\psi|^2$ means that the effective
rest mass varies with the field strength, a self-gravitational
backreaction absent in the flat-background treatment.

The self-consistent system can also be read as a concrete implementation
of the programme suggested by Shojai and
Shojai~\cite{ShojaiShojai1998}: the quantum potential appears as
spacetime curvature, and the transition from quantum to classical
propagation corresponds to geometric relaxation toward the vacuum.
Whether this system admits non-trivial stationary solutions, and
whether the resulting quantisation conditions have independent physical
content, is addressed in the next subsection.
\end{remark}

\subsection{Exact solutions and spectral quantisation}
\label{subsec:sc-solutions}

We now show that the self-consistent system admits non-trivial exact
solutions.  The key is a linearising substitution: the map
$\Phi:=|\psi|^3$ converts the nonlinear field equation into a classical
PDE with well-understood spectral theory.

\begin{theorem}[Cube-harmonic linearisation]
\label{thm:cube-harmonic}
Let $\psi=A(x)\,e^{-i\omega t}$ with $A>0$ time-independent.  Then the
massless ($m=0$) self-consistent equation~\eqref{eq:sc-flat} is equivalent
to the Helmholtz equation for $\Phi:=A^3$:
\begin{equation}
  \label{eq:helmholtz}
  \Delta\Phi + 3\omega^2\,\Phi = 0.
\end{equation}
In the static limit ($\omega=0$) this reduces to $\Delta(A^3)=0$: the
cube of the amplitude is harmonic.
\end{theorem}

\begin{proof}
For $\psi=Ae^{-i\omega t}$ with $m=0$, equation~\eqref{eq:sc-flat}
becomes (after dividing by $e^{-i\omega t}$)
\[
  \Delta A+\frac{2\,|\nn A|^2}{A}+\omega^2 A=0.
\]
The identity
$A^2\Delta A+2A\,|\nn A|^2=\nn\cdot(A^2\nn A)=\tfrac{1}{3}\Delta(A^3)$
lets us multiply through by $A^2$ to obtain~\eqref{eq:helmholtz}.
The continuity equation~\eqref{eq:sc-continuity} is identically satisfied
because $\theta=-\omega t$ and $A$ is time-independent:
$\pa_\mu(A^4\pa^\mu\theta)=-\omega\,\pa_t(A^4)=0$.
\end{proof}

\begin{corollary}[Static exact solution]
\label{cor:static-exact}
In the static massless case ($\omega=0$), the unique spherically symmetric
solution with $A>0$ for $r>0$ is
\begin{equation}
  \label{eq:static-sol}
  A(r) = \bigl(a+b/r\bigr)^{1/3},
  \qquad a>0,\;\;b>0,
\end{equation}
generating the wavefunction metric
$g_{\mu\nu}=(a+b/r)^{2/3}\,\eta_{\mu\nu}$.
The scalar curvature is $R=4b^2/[3\,r^4(a+b/r)^{8/3}]>0$ for all $r>0$.
\end{corollary}

\begin{proof}
$\Delta(A^3)=0$ with spherical symmetry gives $A^3=a+b/r$ (the general
harmonic function on $\mathbb{R}^3\setminus\{0\}$).  Positivity requires
$a,b>0$.  The scalar curvature follows from
Corollary~\ref{cor:R-quantum}:
$R=-6Q/A^2$ with $Q=\Delta A/A=-2|\nn A|^2/A^2$
(using $\Delta A=-2|\nn A|^2/A$ from the original equation with
$\omega=0$), giving $R=12|\nn A|^2/A^4$.  Substituting
$A'=-b/(3r^2)(a+b/r)^{-2/3}$ yields the stated expression.
\end{proof}

\begin{remark}[Physical interpretation of the static solution]
\label{rem:static-physics}
The spacetime $(a+b/r)^{2/3}\eta$ has positive curvature decaying as
$r^{-4}$ at infinity and a curvature singularity at $r=0$ where
$R\sim r^{-4}$.  This singularity is \emph{integrable}: the
Kretschner scalar satisfies
$\int_0^\epsilon R^2\,r^2\,dr<\infty$ for any $\epsilon>0$,
so it is milder than, for instance, the $r^{-6}$ Schwarzschild
singularity.  Nevertheless, the metric degenerates at $r=0$
($g_{\mu\nu}\to\infty$ as $(b/r)^{2/3}$), so $r=0$ is a genuine
geometric singularity, not merely a coordinate artifact.  The conformal
factor $A^2=(a+b/r)^{2/3}$ is smooth for $r>0$ and approaches
$a^{2/3}$ at infinity, so the geometry is asymptotically conformally
flat.

Physically, the field $\psi=A(r)$ generates its own effective
stress-energy and propagates self-consistently on the resulting
geometry.  The parameter $b$ sets the``geometric charge''
(strength of the curvature), while $a$ determines the asymptotic
conformal factor $g\to a^{2/3}\eta$ at large~$r$.
\end{remark}

\begin{theorem}[Spectral quantisation from geometric self-consistency]
\label{thm:spectral-quant}
Let $\omega>0$ and consider regular spherically symmetric solutions of the
massless self-consistent system with $A>0$ on $0<r<R$ and $A(R)=0$
(conformal degeneration at $r=R$).  Then:
\begin{enumerate}[label=\textup{(\roman*)},noitemsep]
\item The frequency is uniquely determined by the cavity radius:
\begin{equation}
  \label{eq:quant-freq}
  \omega = \frac{\pi}{\sqrt{3}\,R}\,.
\end{equation}
\item The amplitude profile is
\begin{equation}
  \label{eq:quant-profile}
  A(r)=\left(\frac{C\,\sin(\pi r/R)}{r}\right)^{\!1/3},
  \qquad C>0.
\end{equation}
\item No higher harmonics ($n\ge 2$) are admissible: the positivity
  constraint $A=\Phi^{1/3}>0$ selects the ground-state eigenfunction.
\end{enumerate}
\end{theorem}

\begin{proof}
By Theorem~\ref{thm:cube-harmonic}, $\Phi=A^3$ satisfies the
Helmholtz equation $\Delta\Phi+3\omega^2\Phi=0$.  Regular spherically
symmetric solutions are
$\Phi=\hat C\,\sin(\sqrt{3}\omega\,r)/r$.  The boundary condition
$\Phi(R)=0$ requires $\sqrt{3}\omega R=n\pi$ for some positive
integer $n$, giving $\omega_n=n\pi/(\sqrt{3}R)$.

For $n\ge 2$, $\Phi$ has interior zeros (at $r=kR/n$, $k=1,\dots,n-1$).
Since $\Phi=A^3$ and $A\ge 0$, these zeros force $A=0$ at interior
points, contradicting $A>0$ on $(0,R)$.  Hence only $n=1$ is admissible,
and $\omega=\pi/(\sqrt{3}R)$.

Equivalently, by the Courant nodal domain theorem the only positive
Dirichlet eigenfunction of $-\Delta$ on a ball is the first
eigenfunction, which is non-degenerate.
\end{proof}

\begin{proposition}[Curvature invariants of the self-consistent cavity]
\label{prop:cavity-curvature}
For the oscillating massless solution of
Theorem~\ref{thm:spectral-quant}, let $\Phi=A^3=C\sin(\pi r/R)/r$.
The scalar curvature of the self-generated metric $g=A^2\eta$ is
\begin{equation}
  \label{eq:cavity-R}
  R_g
  = \frac{6\omega^2}{A^2}
    + \frac{4(\Phi')^2}{3\,\Phi^{8/3}}\,,
  \qquad \omega=\frac{\pi}{\sqrt{3}\,R}\,.
\end{equation}
At the origin, $A(0)=(C\pi/R)^{1/3}$ is finite and $A'(0)=0$ by
symmetry, so
$R_g(0)=6\omega^2/(C\pi/R)^{2/3}<\infty$: the interior geometry is
regular.  At the cavity boundary $r\to R^-$, $A\sim(R-r)^{1/3}\to 0$
and $R_g\sim(R-r)^{-8/3}\to\infty$: the scalar curvature diverges,
consistent with the metric degeneration $g\to 0$
(Corollary~\ref{cor:nodes}).
\end{proposition}

\begin{proof}
From Corollary~\ref{cor:R-quantum}, $R_g=-6\Delta A/A^3$.
Since $A=\Phi^{1/3}$,
$\Delta A=\tfrac{1}{3}\Phi^{-2/3}\Delta\Phi
 -\tfrac{2}{9}\Phi^{-5/3}|\nn\Phi|^2$.
Substituting $\Delta\Phi=-3\omega^2\Phi$ gives
$\Delta A=-\omega^2 A-2|\nn\Phi|^2/(9\Phi^{5/3})$,
and $A^3=\Phi$ yields~\eqref{eq:cavity-R}.
Near $r=R$, $\Phi\sim C\pi(R-r)/R^2$ so
$\Phi^{8/3}\sim(R-r)^{8/3}$ while $(\Phi')^2\to(C\pi/R^2)^2>0$,
giving the exponent $-8/3$.
\end{proof}

\begin{remark}[Self-generated geometric cavity]
\label{rem:cavity}
The ``cavity'' in Theorem~\ref{thm:spectral-quant} is not imposed
externally: the metric $g=A^2\eta$ degenerates at $r=R$ where $A\to 0$
(cf.\ Corollary~\ref{cor:nodes}).  The field creates its own confining
geometry.

Importantly, the quantisation in~\eqref{eq:quant-freq} is \emph{relational}
rather than absolute: for any cavity radius $R>0$ there exists a
self-consistent solution at $\omega=\pi/(\sqrt{3}\,R)$, and conversely
any $\omega>0$ determines a cavity of radius $R=\pi/(\sqrt{3}\,\omega)$.
What is rigid is the $\omega$--$R$ ratio, not the individual values.
Nevertheless, this constraint is non-trivial: it arises purely from
geometric self-consistency (the field propagates on the metric it
generates) together with regularity at the origin and positivity
($|\psi|>0$ in the interior), without any external quantisation
postulate.  If an independent mechanism fixes $R$ (e.g.\  boundary
conditions, a mass scale from the $m>0$ theory, or matching to an
external geometry), the frequency becomes absolutely quantised.
\end{remark}

\begin{proposition}[Massive case: Emden--Fowler reduction]
\label{prop:emden-fowler}
For $m>0$ and the time-periodic ansatz
$\psi=A(r)\,e^{-i\omega t}$, the self-consistent
equation~\eqref{eq:sc-flat} is equivalent to
\begin{equation}
  \label{eq:emden-fowler}
  \Delta\Phi+3\omega^2\,\Phi+3m^2\,\Phi^{5/3}=0,
  \qquad \Phi=A^3.
\end{equation}
This is a nonlinear eigenvalue problem of Emden--Fowler type
(exponent $5/3$).  The mass term $m^2\Phi^{5/3}$ couples the amplitude
to the frequency, breaking the conformal scaling symmetry of the
massless case and potentially selecting isolated solution branches.
\end{proposition}

\begin{proof}
Proceeding as in Theorem~\ref{thm:cube-harmonic} but retaining $m$:
multiplying the radial equation by $A^2$ and using
$\nabla\!\cdot\!(A^2\nabla A)=\tfrac{1}{3}\Delta(A^3)$ gives
$\tfrac{1}{3}\Delta\Phi+\omega^2\Phi+m^2 A^5=0$.  Since
$A^5=(\Phi)^{5/3}$, this is~\eqref{eq:emden-fowler}.
\end{proof}

\begin{remark}[Amplitude--frequency coupling in the massive case]
\label{rem:massive-quant}
In the massless system, the Helmholtz equation~\eqref{eq:helmholtz} is
linear: the amplitude $C$ in~\eqref{eq:quant-profile} is a free parameter
and does not affect the frequency.  For $m>0$, the
nonlinearity~\eqref{eq:emden-fowler} forces $\omega$ to depend on both
the cavity size $R$ and the amplitude $C$.  This
amplitude--frequency coupling is the hallmark of a self-gravitational
backreaction: the field's own energy content shifts its oscillation
frequency.  The following results establish that this coupling uniquely
determines the amplitude once the frequency is fixed, and that an
additional charge constraint produces discrete solution branches.
\end{remark}

\begin{theorem}[Existence, uniqueness, and domain for the massive system]
\label{thm:massive-existence}
Let $m>0$ and $R>0$.  The rescaling
$\Phi(r)=(3m^2)^{-3/2}R^{-3}\,\hat\Phi(r/R)$
transforms~\eqref{eq:emden-fowler} on $B_R$ into
\begin{equation}
  \label{eq:rescaled-ef}
  \Delta\hat\Phi+\Lambda\hat\Phi+\hat\Phi^{5/3}=0,
  \quad \hat\Phi>0\;\text{on }B_1,
  \quad\hat\Phi\big|_{\pa B_1}=0,
\end{equation}
where $\Lambda:=3\omega^2 R^2$.  Then:
\begin{enumerate}[label=\textup{(\roman*)},noitemsep]
\item For each $\Lambda\in[0,\pi^2)$,
  problem~\eqref{eq:rescaled-ef} has exactly one positive radial
  solution~$\hat\Phi_\Lambda$.
\item For $\Lambda\ge\pi^2$, no positive solution exists.
\item The map $\Lambda\mapsto\hat\Phi_\Lambda(0)$ is continuous and
  positive on $[0,\pi^2)$, with
  $\hat\Phi_\Lambda(0)\to 0$ as $\Lambda\nearrow\pi^2$.
\end{enumerate}
\end{theorem}

\begin{proof}
The exponent $5/3$ is subcritical ($5/3<5=(n+2)/(n-2)$ for $n=3$).

\emph{Existence.}  For $\Lambda\in(0,\pi^2)$ the functional
$E(u)=\tfrac{1}{2}\!\int|\nn u|^2
 -\tfrac{\Lambda}{2}\!\int u^2
 -\tfrac{3}{8}\!\int|u|^{8/3}$
has mountain-pass geometry on $H^1_0(B_1)$; the subcritical exponent
guarantees the Palais--Smale condition via the compact embedding
$H^1_0\hookrightarrow L^{8/3}$.  For $\Lambda=0$ the same argument
applies on the Nehari manifold with the pure nonlinearity
$\hat\Phi^{5/3}$.

\emph{Uniqueness.}  By Gidas--Ni--Nirenberg, every positive
solution is radially symmetric and radially decreasing, hence determined
by its central value $\alpha=\hat\Phi(0)$.  Denoting by $R_1(\alpha)$
the first zero of the radial ODE solution with initial data
$\hat\Phi(0)=\alpha$, $\hat\Phi'(0)=0$, Sturm comparison shows that
$R_1(\alpha)$ is strictly monotone in $\alpha$ (larger central values
produce stronger nonlinear restoring and earlier zeroes), so at most
one $\alpha$ satisfies $R_1(\alpha)=1$.

\emph{Non-existence.}  Testing~\eqref{eq:rescaled-ef} against
the first eigenfunction $\phi_1>0$ of $-\Delta$ on $B_1$
(eigenvalue $\lambda_1=\pi^2$):
$\int(\pi^2-\Lambda)\hat\Phi\,\phi_1=\int\hat\Phi^{5/3}\phi_1>0$,
which requires $\Lambda<\pi^2$.

\emph{Continuity.}  The solution branch bifurcates (supercritically)
from the trivial solution at $\Lambda=\pi^2$: near the bifurcation,
$\hat\Phi_\Lambda\approx\epsilon\,\phi_1$ with
$\epsilon\sim C(\pi^2-\Lambda)^{3/2}\to 0$.  The branch extends
continuously to $\Lambda=0$ by the implicit function theorem
(the linearised operator is invertible for each $\Lambda<\pi^2$
since the positive solution is unique).
\end{proof}

\begin{corollary}[Amplitude--frequency locking and gravitational redshift]
\label{cor:amp-freq-lock}
Fix $m>0$ and $R>0$, and let
$\omega_0:=\pi/(\sqrt{3}\,R)$ denote the massless cut-off frequency.
\begin{enumerate}[label=\textup{(\roman*)},noitemsep]
\item For each $\omega\in(0,\omega_0)$, the self-consistent massive
  system has a unique positive radial solution whose central amplitude
  $A(0)=\Phi(0)^{1/3}$ is uniquely determined by $\omega$.
\item The mass term always shifts the frequency below the massless
  value: $\omega<\omega_0$.  Larger amplitudes correspond to lower
  frequencies---a gravitational redshift arising from the field's own
  conformal factor.
\item As $\omega\nearrow\omega_0$, the amplitude vanishes and the
  solution reduces to the linearised Helmholtz mode of
  Theorem~\ref{thm:spectral-quant} (rescaled to zero amplitude).
  As $\omega\searrow 0$, the solution approaches the massive static
  configuration $\Delta\Phi+3m^2\Phi^{5/3}=0$
  (the Lane--Emden limit).
\end{enumerate}
\end{corollary}

\begin{proof}
$\omega\in(0,\omega_0)$ corresponds to
$\Lambda=3\omega^2R^2\in(0,\pi^2)$.  By
Theorem~\ref{thm:massive-existence}, the unique solution
$\hat\Phi_\Lambda$ determines the physical amplitude
$\Phi(0)=(3m^2)^{-3/2}R^{-3}\hat\Phi_\Lambda(0)$.
Near the bifurcation ($\Lambda\nearrow\pi^2$), the nonlinear term
becomes negligible and the solution approaches a rescaled Helmholtz
eigenfunction; at $\Lambda=0$ the linear term drops out and only the
Lane--Emden nonlinearity survives.
\end{proof}

\begin{proposition}[Charge quantisation from the fold structure]
\label{prop:charge-quant}
The continuity equation~\eqref{eq:sc-continuity} with
$\theta=-\omega t$ defines a conserved geometric charge
\begin{equation}
  \label{eq:charge}
  \mathcal{Q}
  :=4\pi\omega\!\int_0^R\! \Phi^{4/3}\,r^2\,dr,
  \qquad \Phi=A^3.
\end{equation}
For fixed $m>0$ and $R>0$:
\begin{enumerate}[label=\textup{(\roman*)},noitemsep]
\item $\mathcal{Q}$ is a continuous function of
  $\Lambda=3\omega^2R^2$ on $(0,\pi^2)$ with
  $\mathcal{Q}\to 0$ at both endpoints.
\item There exists a maximum charge
  $\mathcal{Q}_{\max}(m,R)>0$, and no self-consistent
  configuration of radius~$R$ exists for
  $\mathcal{Q}>\mathcal{Q}_{\max}$.
\item For each $\mathcal{Q}\in(0,\mathcal{Q}_{\max})$, the equation
  $\mathcal{Q}(\Lambda)=\mathcal{Q}$ has at least two solutions:
  a \emph{weakly gravitating mode}
  (high~$\omega$, low amplitude, near the Helmholtz limit) and a
  \emph{strongly gravitating mode}
  (low~$\omega$, high amplitude, near the Lane--Emden limit).
\end{enumerate}
\end{proposition}

\begin{proof}
In rescaled variables (Theorem~\ref{thm:massive-existence}),
\[
  \mathcal{Q}(\Lambda)
  =\frac{4\pi\sqrt{\Lambda/3}}{(3m^2)^2 R^2}\,I(\Lambda),
  \qquad
  I(\Lambda):=\int_0^1\!\hat\Phi_\Lambda^{4/3}\,\sigma^2\,d\sigma.
\]
As $\Lambda\nearrow\pi^2$:
$\hat\Phi_\Lambda\to 0$ uniformly, so $I(\Lambda)\to 0$
and $\mathcal{Q}\to 0$.
As $\Lambda\searrow 0$:
$I(\Lambda)\to I(0)>0$ (the Lane--Emden solution has finite
$L^{4/3}$ norm) but $\sqrt\Lambda\to 0$, so $\mathcal{Q}\to 0$.
Extending $\mathcal{Q}$ to $[0,\pi^2]$ by continuity, the maximum
is attained on this compact interval at some interior point.
For $\mathcal{Q}<\mathcal{Q}_{\max}$, the intermediate
value theorem gives at least two pre-images (one on each side
of the maximum).
\end{proof}

\begin{remark}[Physical interpretation of the fold]
\label{rem:fold-physics}
The weakly gravitating branch is a near-linear perturbation of the
massless Helmholtz eigenstate: the backreaction is a small correction.
The strongly gravitating branch is a genuinely nonlinear configuration
in which the conformal factor $A^2$ varies significantly and the
geometry departs from flat.

The maximum charge scales as
$\mathcal{Q}_{\max}\propto 1/(m^4 R^2)$, providing
an upper bound on the charge supportable by a self-consistent
oscillating configuration of given radius---a self-gravitational
analogue of a Chandrasekhar-type limit.

In the massless system, the charge
$\mathcal{Q}=4\pi\omega C^{4/3}\int_0^R(\sin(\pi r/R)/r)^{4/3}r^2\,dr$
depends independently on $\omega$ (which is locked to $R$) and the
free amplitude $C$, so any $\mathcal{Q}>0$ is achievable by
adjusting~$C$.  The mass introduces the coupling between $\omega$
and $C$ that produces the fold, and it is this coupling that
elevates the relational $\omega$--$R$ quantisation
of Theorem~\ref{thm:spectral-quant} to the amplitude quantisation
described here.
\end{remark}

% ================================================================
\section{Extension to Complex Seams}
\label{sec:complex}
% ================================================================

Remark~\ref{rem:complex-seam} identified three structural limitations of
the real scalar seam: vacuum flatness rigidity, a single functional degree
of freedom, and identically vanishing vorticity.  In this section we promote
the seam to a complex-valued field $\varphi=u+iv$ and verify that each
limitation is resolved.

\subsection{The complex seam metric}

\begin{definition}[Complex seam metric on a real background]
\label{def:complex-seam-metric}
Let $(M^n,h)$ be a (pseudo-)Riemannian manifold and
$\varphi = u + iv \colon M \to \mathbb{C}$ a smooth complex-valued
function.  A \emph{complex seam metric} of order~$\le 2$ is
\begin{equation}
\label{eq:complex-seam-metric}
g_{\mu\nu}
= \alpha\,h_{\mu\nu}
  + \beta_u\,\pa_\mu u\,\pa_\nu u
  + \beta_v\,\pa_\mu v\,\pa_\nu v
  + \beta_\times\,(\pa_\mu u\,\pa_\nu v + \pa_\nu u\,\pa_\mu v)
  + \gamma_u\,\nn_\mu\nn_\nu u
  + \gamma_v\,\nn_\mu\nn_\nu v,
\end{equation}
where $\alpha,\beta_u,\beta_v,\beta_\times,\gamma_u,\gamma_v$ are
smooth real-valued functions of the scalar jet invariants
$(u,v,|\nn u|^2_h,|\nn v|^2_h,
\langle\nn u,\nn v\rangle_h,\Delta_h u,\Delta_h v)$.
\end{definition}

\begin{remark}[Recovery of the real ansatz]
\label{rem:recovery-real}
Setting $v \equiv 0$ reduces~\eqref{eq:complex-seam-metric} to the
real seam metric~\eqref{eq:seam-metric} (with $\beta_v$, $\beta_\times$,
and $\gamma_v$ absent and $s = u$).  The real framework is a sub-family
of the complex one.
\end{remark}

\subsection{Two degrees of freedom}

\begin{proposition}[Functional independence]
\label{prop:complex-dof}
Let $\varphi = u + iv$ with $u,v \in C^\infty(M)$ independent
(i.e.\ $du \wedge dv \neq 0$ on a dense open set).
The conformal-plus-gradient metric
\begin{equation}
  \label{eq:conf-grad-complex}
  g_{\mu\nu}
  = e^{2u}\,\eta_{\mu\nu}
    + \beta\,\pa_\mu v\,\pa_\nu v
\end{equation}
(with $\beta = \beta(u,v)$ not identically zero) defines a two-function
family of non-isometric metrics.  At a point where $du \wedge dv \neq 0$,
the metric has eigenvalues in $n-2$ directions controlled by $e^{2u}$
and distinct eigenvalues in the remaining two directions controlled
independently by $u$ and $\beta$, giving two spectral degrees of
freedom compared with one in the real case.
\end{proposition}

\begin{proof}
Write the metric as $g = e^{2u}\eta + \beta\,dv \otimes dv$.
In an $\eta$-orthonormal frame with $e_1 \propto \nn u$ and
$e_2 \propto \nn v - \frac{\langle\nn v,\nn u\rangle}{|\nn u|^2}\nn u$,
the $n-2$ directions orthogonal to both gradients see eigenvalue
$e^{2u}$, while the $e_1$- and $e_2$-directions have distinct eigenvalues
modified by $\beta$.  Since $u$ and $v$ are independent, the two
spectral parameters are independently adjustable; in particular, two
choices $(u_1,v_1) \neq (u_2,v_2)$ with $du_i \wedge dv_i \neq 0$
give metrics with generically distinct scalar curvatures, hence are
non-isometric.
\end{proof}

\subsection{Vorticity from the cross-gradient}

In the real-seam framework, whenever $\nn s$ is timelike the unit
4-velocity $U = \nn s / |\nn s|$ has vanishing vorticity because $U$
is proportional to the exact 1-form $ds$
(Corollary~\ref{cor:irrotational}).  The complex ansatz breaks this.

\begin{theorem}[Non-vanishing vorticity from a complex seam]
\label{thm:complex-vorticity}
Let $g$ be a complex seam metric~\eqref{eq:complex-seam-metric} on
Minkowski background in $n = 4$.  Define the 1-form
\begin{equation}
\label{eq:tilted-velocity}
W_\mu := \pa_\mu u + \lambda(x)\,\pa_\mu v,
\end{equation}
where $\lambda$ is a smooth function.  Then the vorticity 2-form satisfies
\begin{equation}
\label{eq:vorticity-formula}
dW = d\lambda \wedge dv.
\end{equation}
In particular, if $\lambda = \lambda(u)$, then
$dW = \lambda'(u)\,du \wedge dv \neq 0$
wherever $\lambda'(u) \neq 0$ and $du \wedge dv \neq 0$.
In the real limit $v \to 0$ the support vanishes, recovering the
irrotationality of the real seam.
\end{theorem}

\begin{proof}
Since $du$ is exact, $d(du) = 0$.  Hence
$dW = d(\lambda\,dv) = d\lambda \wedge dv$.
If $\lambda = \lambda(u)$, then $d\lambda = \lambda'\,du$, giving
$dW = \lambda'\,du \wedge dv$.
\end{proof}

\begin{remark}[Physical interpretation]
\label{rem:vorticity-physics}
The 1-form $W$ is not the normalised 4-velocity of the seam flow, but
the theorem demonstrates that the complex ansatz provides a tensorial
mechanism for vorticity \emph{absent} in the real case: the 2-form
$du \wedge dv$ is identically available once $\varphi$ is complex.
In a charged scalar-field context, the macroscopic current
$J^\mu \propto |\psi|^2\,\nn^\mu\theta$ naturally takes a form
involving both $\nn u$ and $\nn v = \nn\theta$; any physical
4-velocity built from both gradients will generically inherit nonzero
vorticity.
\end{remark}

\subsection{Breaking vacuum flatness rigidity}

A natural question is whether the complex ansatz can break the vacuum
flatness rigidity of the real seam
(Theorems~\ref{thm:vacuumon}, \ref{thm:vacuum};
Corollary~\ref{cor:lorentz-vacuum}).
Within the conformal-plus-phase-gradient subfamily
\begin{equation}
\label{eq:conf-phase-metric}
g_{\mu\nu} = e^{2u}\bigl(\eta_{\mu\nu} + \beta\,\pa_\mu v\,\pa_\nu v\bigr),
\end{equation}
rigidity persists: for static $u = -c\log|z|$ and linear
$v = m\,x^0$, the metric is conformally flat and a direct computation
gives $G_{33} = 3c^2/z^2$, forcing $c = 0$.
However, the full six-generator span admits genuinely curved vacua.

\begin{theorem}[Non-flat vacuum from a complex seam]
\label{thm:complex-vacuum}
Let $u = \log z$, $v = m\,t$ (with $m > 0$, $z > 0$), and set
\begin{equation}
\label{eq:taub-coefficients}
\alpha(u) = e^{4u/3},\qquad
\beta_v(u) = \frac{e^{4u/3} - e^{-2u/3}}{m^2},\qquad
\gamma_u(u) = e^{10u/3} - e^{2u},
\end{equation}
with $\beta_u = \beta_\times = \gamma_v = 0$.  Then the complex seam
metric~\eqref{eq:complex-seam-metric} on Minkowski background
$(M,\eta)$ is
\begin{equation}
\label{eq:taub-metric}
g = \mathrm{diag}\!\bigl(-z^{-2/3},\;z^{4/3},\;z^{4/3},\;1\bigr),
\end{equation}
the Taub plane-symmetric vacuum.  This metric is Ricci-flat
($G_{\mu\nu} = 0$) and genuinely curved: the Riemann tensor has
$R^0_{\;101} = \tfrac{2}{9}\,z^{-2/3} \neq 0$.
\end{theorem}

\begin{proof}
Since $u = \log z$ we have $\pa_\mu u = \delta_{\mu 3}/z$ and
$\nn_\mu\nn_\nu u = -\delta_{\mu 3}\delta_{\nu 3}/z^2$;
since $v = m\,t$ we have $\pa_\mu v = m\,\delta_{\mu 0}$.
The six generators evaluate as follows.
\begin{enumerate}[label=\quad(\arabic*)]
\item $\alpha\,\eta_{\mu\nu}$: contributes
  $-z^{4/3}$ to $g_{00}$,\ $z^{4/3}$ to $g_{11}$ and $g_{22}$,\
  $z^{4/3}$ to $g_{33}$.
\item $\beta_v\,\pa_\mu v\,\pa_\nu v = (z^{4/3}-z^{-2/3})\,\delta_{\mu 0}\delta_{\nu 0}$:
  adds $z^{4/3}-z^{-2/3}$ to $g_{00}$.
\item $\gamma_u\,\nn_\mu\nn_\nu u = -(z^{10/3}-z^2)/z^2\,\delta_{\mu 3}\delta_{\nu 3}$:
  adds $1 - z^{4/3}$ to $g_{33}$.
\end{enumerate}
Summing: $g_{00} = -z^{4/3} + z^{4/3} - z^{-2/3} = -z^{-2/3}$,\
$g_{11} = g_{22} = z^{4/3}$,\ $g_{33} = z^{4/3} + 1 - z^{4/3} = 1$.
This is~\eqref{eq:taub-metric}.  The Ricci-flat and Riemann
statements are verified by direct computation (the metric depends on
$z$ alone, so all curvature reduces to ODEs in $z$).
\end{proof}

\begin{remark}[Mechanism and real-seam obstruction]
\label{rem:rigidity-persists}
In the real seam the conformal generator forces
$g_{00}/g_{11} = -1$, but the Taub metric requires
$g_{00}/g_{11} = -z^{-2}$.  The phase-gradient generator
$\beta_v\,\pa_\mu v\,\pa_\nu v$ breaks this lock by contributing
only to $g_{00}$, independently of the conformal factor;
the Hessian generator $\gamma_u\,\nn_\mu\nn_\nu u$ similarly
adjusts $g_{33}$ alone.  Together the three generators $\alpha$,
$\beta_v$, $\gamma_u$ supply three independent diagonal degrees of
freedom, precisely what is needed to realise an anisotropic vacuum.
The conformal-plus-gradient subfamily~\eqref{eq:conf-phase-metric}
lacks both $\beta_v$ and $\gamma_u$ as functions of $u$
independently of $\alpha$, which is why it inherits the
real-seam rigidity.
\end{remark}

\subsection{Self-consistent complex wavefunction seam}

The wavefunction seam of Section~\ref{sec:wavefunction} identifies
the seam scalar with the log-amplitude of a Klein--Gordon field:
$s = \log|\psi|$, $\psi = |\psi|\,e^{i\theta}$.  The complex seam
provides the natural extension: setting $u = \log|\psi|$ and $v = \theta$
promotes \emph{both} quantum-mechanical degrees of freedom---amplitude and
phase---to geometric data.

\begin{definition}[Complex wavefunction seam metric]
\label{def:complex-wf-seam}
Let $\psi = |\psi|\,e^{i\theta}$ be a smooth complex scalar field with
$|\psi| > 0$.  The \emph{complex wavefunction seam metric} is the complex
seam metric~\eqref{eq:complex-seam-metric} with
$u = \log|\psi|$ and $v = \theta$.
The coefficient functions depend on the scalar jet invariants
of $(\log|\psi|,\theta)$.
Setting $\beta_v = \beta_\times = \gamma_u = \gamma_v = 0$ and
$\alpha = |\psi|^2$ recovers the conformal wavefunction seam
$g = |\psi|^2\eta$ of Section~\ref{sec:wavefunction}.
\end{definition}

\begin{proposition}[Vacuum no-go: conformal vs.\ complex]
\label{prop:vacuum-nogo-lift}
\mbox{}
\begin{enumerate}[label=\upshape(\roman*)]
\item In the conformal wavefunction seam $g = |\psi|^2\eta$, every
Ricci-flat metric is flat (Theorem~\ref{thm:vacuum}).
\item In the complex wavefunction seam
(Definition~\ref{def:complex-wf-seam}),
Ricci-flat metrics need not be flat: the wavefunction
$\psi = z\,e^{imt}$ generates the Taub vacuum~\eqref{eq:taub-metric}
via the coefficient functions of Theorem~\ref{thm:complex-vacuum}.
\end{enumerate}
\end{proposition}

\begin{definition}[Self-consistent complex wavefunction system]
\label{def:sc-complex}
A \emph{self-consistent complex wavefunction seam} is a pair $(\psi,g)$
where $g$ is the complex wavefunction seam
metric (Definition~\ref{def:complex-wf-seam}) and $\psi$ satisfies the
Klein--Gordon equation on the spacetime it generates:
\begin{equation}
\label{eq:sc-complex-kg}
(\Box_g + m^2)\,\psi = 0.
\end{equation}
The conformal restriction $\alpha = |\psi|^2$,
$\beta_u = \beta_v = \beta_\times = \gamma_u = \gamma_v = 0$ recovers the
self-consistent system of
Section~\ref{subsec:self-consistent}, whose solutions
exhibit spectral quantisation (Theorem~\ref{thm:spectral-quant}).
\end{definition}

\begin{remark}[The phase as a geometric degree of freedom]
\label{rem:phase-geometry}
In the conformal wavefunction seam, the phase $\theta$ is geometrically
invisible: the metric $g = |\psi|^2\eta$ depends only on the amplitude.
The phase enters the physics only through the Klein--Gordon constraint,
which in the vacuum limit reduces to the eikonal equation
$|\nn\theta|^2 = m^2$ (Theorem~\ref{thm:kg-compat})---a condition
\emph{on} $\theta$ but not \emph{from} $\theta$ in the metric.

In the complex wavefunction seam, the phase enters the metric directly
through the generators $\beta_v\,\pa_\mu\theta\,\pa_\nu\theta$
and $\gamma_v\,\nn_\mu\nn_\nu\theta$.  It is this promotion
of the phase from constraint-level to metric-level that
breaks the conformal rigidity
(Theorem~\ref{thm:complex-vacuum},
Remark~\ref{rem:rigidity-persists}):
the phase-gradient generator decouples the temporal and spatial
conformal factors, enabling the anisotropic metric structure required
by non-flat vacuum solutions.
In this sense, the phase of the wavefunction---a second degree of
freedom absent from the real seam and supplied by the
quantum-mechanical structure of $\psi$---is the geometric ingredient
that enables curvature in the self-consistent framework.
\end{remark}

% ================================================================
\section{Conclusion}
\label{sec:conclusion}
% ================================================================

We have proved eight structural results.

For the conformally flat metric $g_{\mu\nu}=e^{2s}\eta_{\mu\nu}$,
the Einstein tensor is an exact decomposition into three seam generators
(Theorem~\ref{thm:conformal}), with anisotropic tidal stress, kinetic
flux, and isotropic pressure corresponding respectively to the Hessian,
gradient-tensor, and conformal generators of the seam decomposition.
The vacuum condition reduces to harmonicity of $e^s$
(Theorem~\ref{thm:vacuum}).

For the $O(n)$-invariant family $g_{ij}=A(r)\delta_{ij}+B(r)x_ix_j$,
the Einstein tensor takes the Cartesian form $C_1\delta_{ij}+C_2 x_ix_j$
with explicit coefficients (Theorem~\ref{thm:on}).  No generator beyond
the two surviving seam generators appears.  The vacuum condition is
rigid: it forces the metric to be flat within this family
(Theorem~\ref{thm:vacuumon}).

For the Lorentz-invariant family $g_{\mu\nu}=A(s)\eta_{\mu\nu}+B(s)x_\mu x_\nu$,
the metric is equivalent (inside the light cone) to an open ($k=-1$) FLRW
ansatz with explicit lapse and scale factor (Section~\ref{sec:lorentz}).  The
Einstein tensor is computed in that form (Theorem~\ref{thm:lorentz}), and the
vacuum condition forces flatness within this family
(Corollary~\ref{cor:lorentz-vacuum}).

For the wavefunction seam $s=\log|\psi|$ with $\psi$ satisfying the
Klein--Gordon equation, the seam-geometric vacuum condition is equivalent to
the eikonal equation $|\nabla\theta|^2_\eta=m^2$ for the phase---equivalently,
to the vanishing of the relativistic quantum potential
(Theorem~\ref{thm:kg-compat}).  The null energy condition reduces to
null-convexity of $1/|\psi|$ (Proposition~\ref{prop:wf-dictionary}).

In these families, metric signature is an output of the seam structure.
A \emph{sphericity} statement for degeneration loci follows from
$O(n)$-equivariance of the seam rule without the field equations
(Corollary~\ref{cor:birkhoff}).

\medskip
\noindent\textbf{No-go interpretation of vacuum flatness.}
The repeated conclusion ``vacuum forces flatness'' should be read as a
rigidity/no-go statement about the seam-generated ansatz classes, rather than
as a deficiency in solution-finding.  Within the $O(n)$-invariant and
Lorentz-invariant families generated from a scalar seam by an order-$\le 2$
isotropic rule, the Einstein equations overdetermine the remaining degrees of
freedom and eliminate nontrivial vacuum curvature.  In both cases this matches
the familiar rigidity of the underlying symmetry-reduced geometries: the
spatially spherically symmetric Riemannian setting on one hand, and the open
FLRW vacuum setting on the other.

The natural continuation is: (a) extend the computation to the full
three-parameter seam metric on a curved background, to test whether the
extended closure property holds in general; and
(b) identify which coefficient choices for $(A,B)$ correspond to specific
physically known spacetimes (Kerr, de~Sitter, Reissner--Nordstr\"{o}m),
thereby completing the dictionary between seam parameters and GR solutions.

\medskip
\noindent\textbf{First-order closure of the seam decomposition.}
The seam decomposition employs four generators
(Definition~\ref{def:seam-decomposition}).  In the three symmetric
families of Results~1--3, the Hessian-square generator collapses into
the span of the other three
(Remark~\ref{rem:why-closure-held}), explaining why three generators
suffice there.  For a general seam whose Hessian has three or more
distinct eigenvalues, all four generators are needed: the
linearised Einstein tensor of the gradient-tensor perturbation
produces an independent Hessian-squared contribution
(Theorem~\ref{thm:closure-fails}).  With all four generators present,
first-order closure holds
(Theorem~\ref{thm:extended-closure}).  Whether closure persists at the
nonlinear level within an extended equivariant ring remains open.

\medskip
\noindent\textbf{What the results provide.}
Taken together, the structural results serve eight purposes.
First, a \emph{translation layer}: energy conditions become scalar
inequalities on the reciprocal seam factor and its Hessian
(Proposition~\ref{prop:ec-dictionary}, Corollary~\ref{cor:ec-temporal}),
replacing full tensor checks with convexity conditions on a single function.
Second, \emph{no-go information}: the rigidity/vacuum-flatness theorems
(Theorems~\ref{thm:vacuumon}, \ref{thm:lorentz},
Corollaries~\ref{cor:lorentz-vacuum}, \ref{cor:lambda-on})
delineate which seam-generated ansatz classes cannot support nontrivial
vacuum curvature, preventing misdirected searches for solutions within them.
Third, a \emph{program for controlled metric construction}: the closure
analysis identifies the equivariant span needed at each order, and the
temporal reductions (Corollary~\ref{cor:ec-temporal}) show where the full
tensor problem collapses to tractable ODE or convexity constraints.
Fourth, a \emph{bridge to quantum mechanics}: the wavefunction seam
(Section~\ref{sec:wavefunction}) shows that the conformal seam framework
naturally absorbs Klein--Gordon fields, with the vacuum condition
corresponding to the classical (WKB/eikonal) limit and the energy
conditions translating into convexity of $1/|\psi|$.
Fifth, a \emph{self-consistent closure}: requiring $\psi$ to satisfy
Klein--Gordon on $g=|\psi|^2\eta$ rather than on $\eta$ yields a
well-defined nonlinear system
(Proposition~\ref{prop:sc-madelung}) whose Madelung form exhibits
an amplitude-dependent mass-shell relation and a geometrically natural
conserved current.  The vacuum fixed points of this system
(Corollary~\ref{cor:sc-vacuum}) satisfy a modified eikonal equation
$|\nn\theta|^2=m^2|\psi|^2+2|\nn\log|\psi||^2$, providing a concrete
arena for studying self-gravitating quantum fields within the seam
framework.
Sixth, \emph{exact solutions and spectral quantisation}: for $m=0$ the
self-consistent system linearises via the cube-harmonic substitution
$\Phi=|\psi|^3$ (Theorem~\ref{thm:cube-harmonic}), producing the static
exact solution $|\psi|=(a+b/r)^{1/3}$
(Corollary~\ref{cor:static-exact}) and, for the time-periodic ansatz, a
Helmholtz eigenvalue problem whose positivity constraint selects the
ground-state eigenfunction and quantises the frequency:
$\omega=\pi/(\sqrt{3}\,R)$
(Theorem~\ref{thm:spectral-quant}).  The confining ``cavity'' is
self-generated by the metric degeneration at $|\psi|=0$
(Remark~\ref{rem:cavity}).  For $m>0$ the system reduces to an
Emden--Fowler equation (Proposition~\ref{prop:emden-fowler}) whose
subcritical nonlinearity admits a unique positive solution for each
admissible frequency $\omega<\pi/(\sqrt{3}\,R)$
(Theorem~\ref{thm:massive-existence}), uniquely locking the amplitude
to the frequency (Corollary~\ref{cor:amp-freq-lock}).  The conserved
charge from the continuity equation vanishes at both frequency
endpoints, producing a fold: for each charge below a maximum
$\mathcal{Q}_{\max}\propto 1/(m^4 R^2)$, there exist (at least) two
distinct modes---a weakly and a strongly gravitating
branch (Proposition~\ref{prop:charge-quant}).  This constitutes genuine
amplitude quantisation arising from geometric self-consistency.
Seventh, a \emph{complete resolution of real-seam limitations}: the
complex seam extension (Section~\ref{sec:complex}) resolves each of the
three structural limitations of the real case by promoting the seam to
$\varphi=u+iv$.  The phase gradient $\nn v$ provides a second spectral
parameter (Proposition~\ref{prop:complex-dof}), the 2-form
$du\wedge dv$ carries non-vanishing vorticity
(Theorem~\ref{thm:complex-vorticity}), and the full six-generator span
admits genuinely curved vacua: the Taub plane-symmetric vacuum is
realised as a complex seam metric with explicit coefficient functions
(Theorem~\ref{thm:complex-vacuum}).
Eighth, a \emph{self-consistent framework for curved quantum geometry}:
the complex wavefunction seam
(Definition~\ref{def:complex-wf-seam}) unifies the quantum bridge of
Results~4--6 with the complex extension of Result~7 by promoting both
degrees of freedom of $\psi = |\psi|\,e^{i\theta}$ to geometric data.
The conformal wavefunction seam is blind to the phase and forces
vacuum flatness; the complex wavefunction seam makes the phase a
geometric degree of freedom that breaks the conformal rigidity
(Proposition~\ref{prop:vacuum-nogo-lift}).  The self-consistent system
(Definition~\ref{def:sc-complex})---Klein--Gordon on the metric the
field generates---defines a nonlinear framework whose solutions
realise quantum-field-generated curved spacetimes;
the conformal restriction recovers the spectrally quantised system
of Results~5--6.

\backmatter

\section*{Acknowledgements}

The author thanks colleagues for discussions on seam geometry and
mathematical physics.  The author discloses the use of AI-assisted
tools to help with drafting, editing, symbolic computation, and
\LaTeX{} code; all mathematical claims, calculations, and final wording
were reviewed and validated by the author.  AI tools are not listed
as authors.

\section*{Declarations}

\begin{itemize}
\item \textbf{Funding.} No funding was received for conducting this study.
\item \textbf{Competing Interests.} The author has no competing interests to declare that are relevant to the content of this article.
\item \textbf{Data Availability.} Data sharing is not applicable to this article as no datasets were generated or analysed during the current study.
\end{itemize}

\begin{thebibliography}{99}

\bibitem{Wald1984}
Wald, R.M.: General Relativity.
University of Chicago Press, Chicago (1984)

\bibitem{MisnerThorneWheeler1973}
Misner, C.W., Thorne, K.S., Wheeler, J.A.: Gravitation.
W.H. Freeman, San Francisco (1973)

\bibitem{Carroll2004}
Carroll, S.M.: Spacetime and Geometry: An Introduction to General Relativity.
Addison-Wesley, San Francisco (2004)

\bibitem{Weinberg1972}
Weinberg, S.: Gravitation and Cosmology: Principles and Applications of the General Theory of Relativity.
Wiley, New York (1972)

\bibitem{HawkingEllis1973}
Hawking, S.W., Ellis, G.F.R.: The Large Scale Structure of Space-Time.
Cambridge University Press, Cambridge (1973)

\bibitem{Bekenstein1993}
Bekenstein, J.D.: The relation between physical and gravitational geometry.
Phys. Rev. D \textbf{48}(6), 3641--3647 (1993)

\bibitem{CliftonFerreiraPadillaSkordis2012}
Clifton, T., Ferreira, P.G., Padilla, A., Skordis, C.: Modified gravity and cosmology.
Phys. Rep. \textbf{513}(1), 1--189 (2012)

\bibitem{Birkhoff1923}
Birkhoff, G.D.: Relativity and Modern Physics.
Harvard University Press, Cambridge, MA (1923)

\bibitem{Wheeler1990}
Wheeler, J.A.: Information, physics, quantum: The search for links.
In: Zurek, W.H. (ed.) Complexity, Entropy, and the Physics
of Information, pp.~354--368. Addison-Wesley, Redwood City, CA (1990)

\bibitem{ShojaiShojai1998}
Shojai, F., Shojai, A.: On the geometrization of Bohmian mechanics: A new approach to
quantum gravity.
Int. J. Mod. Phys. A \textbf{13}(04), 677--693 (1998)

\bibitem{Chowdhury2021}
Chowdhury, S.R., Gupta, K., Kar, S.: Disformal transformations and the motion of a particle in
semi-classical gravity.
Eur. Phys. J. C \textbf{81}(10), 946 (2021)

\bibitem{MorrisThorne1988}
Morris, M.S., Thorne, K.S.: Wormholes in spacetime and their use for interstellar travel:
A tool for teaching general relativity.
Am. J. Phys. \textbf{56}(5), 395--412 (1988)

\bibitem{KMS1993}
Kol\'a\v{r}, I., Michor, P.W., Slov\'ak, J.: Natural Operations in Differential Geometry.
Springer-Verlag, Berlin (1993)

\end{thebibliography}

\end{document}
